\documentclass[lang=cn,11pt]{elegantbook}

\title{概率论与数理统计}
\subtitle{学习笔记}

\author{夏同}
\institute{华南理工大学}
\date{\today}
\version{1.0}
\bioinfo{自定义}{信息}

\extrainfo{居敬持志，守正出奇。—— 张红兵}

\setcounter{tocdepth}{3}

\logo{logo-blue.png}
\cover{cover.jpg}

% 本文档命令
\usepackage{array}
\newcommand{\ccr}[1]{\makecell{{\color{#1}\rule{1cm}{1cm}}}}

% 修改标题页的橙色带
% \definecolor{customcolor}{RGB}{32,178,170}
% \colorlet{coverlinecolor}{customcolor}

\begin{document}

\maketitle
\frontmatter

\tableofcontents

\mainmatter

\chapter{随机事件与概率}
\[
\begin{cases}
\text{基本概念} & 
\begin{cases}
\text{随机试验} \ E: \text{可重复性、结果不唯一、所有结果已知、不确定性} \\
\text{样本空间} \ S: \text{所有可能结果构成的集合} \\
\text{随机事件} \ A: \text{样本空间的子集} 
\begin{cases}
\text{关系：包含、相等、互斥、对立} \\
\text{表示：和事件、积事件、差事件、逆(对立)事件} \\
\text{运算：交换律、结合律、分配律、德摩根根律}
\end{cases}
\end{cases} \\
\text{概率概念} & 
\begin{cases}
\text{统计概念：频率} \\
\text{公理化概念：非负性、规范性和可列可加性} \\
\text{古典模型(等可能概型)：有限元素、可能性相同} \\
\text{计算公式}
\begin{cases}
\text{加法公式}
\begin{cases}
P(A \cup B) = P(A) + P(B) - P(AB) \\
P(A \cup B) = P(A) + P(B) \ (\text{若事件} \ A \ \text{与} \ B \ \text{互斥})
\end{cases} \\
\text{减法公式}
\begin{cases}
P(A - B) = P(A) - P(AB) = P(A\overline{B}) \\
P(A - B) = P(A) - P(B) \ (\text{若事件} \ A \supset B)
\end{cases} \\
\text{求逆公式：} \ P(\overline{A}) = 1 - P(A)
\end{cases}
\end{cases} \\
\text{条件概率概念} & 
\begin{cases}
\text{定义：当} \ P(A) > 0 \ \text{时}, \ P(B|A) = \dfrac{P(AB)}{P(A)}, \text{且满足概率性质与上述公式} \\
\quad (\text{非负性、规范性和可列可加性}) \\
\text{计算:定义法，缩小样本空间方法} \\
\text{应用}
\begin{cases}
\text{乘法公式：当} \ P(A) > 0 \ \text{时}, \ P(AB) = P(B|A)P(A) \\
\text{全概率公式(由因求果)：} \ P(A) = \sum_{i=1}^{n} P(A|B_i)P(B_i) \\
\text{贝叶斯公式(执果寻因)：} \ P(B_i|A) = \dfrac{P(A|B_i)P(B_i)}{\sum_{j=1}^{n} P(A|B_j)P(B_j)} \\
\quad (\text{其中} \ B_1,B_2,\dots,B_n \ \text{是一组完备事件组}, \ P(B_i) > 0)
\end{cases}
\end{cases} \\
\text{独立性} & 
\begin{cases}
\text{定义}
\begin{cases}
\text{两个事件:} \ P(AB) = P(A)P(B) \Leftrightarrow \text{当} \ 0 < P(A) < 1, P(B|A) = P(B|\overline{A}) \\
\text{多个事件：相互独立} \Rightarrow \text{两两独立}, \ \text{两两独立} \nRightarrow \text{相互独立}
\end{cases} \\
\text{性质}
\begin{cases}
\text{相互独立的事件中部分事件间相互独立} \\
\text{相互独立的事件中部分事件的逆事件与其余事件仍相互独立}
\end{cases} \\
\text{应用}
\begin{cases}
\text{计算多个事件的和事件概率转化为逆事件概率乘积} \\
\text{伯努利试验：只有两个结果} \ A \ \text{与} \ \overline{A},  n \ \text{次试验相互独立}
\end{cases}
\end{cases}
\end{cases}
\]
\newpage

\begin{introduction}
  \item 随机试验与样本空间~\ref{def:random-experiment}
  \item 随机事件及其运算~\ref{def:random-event}
  \item 完备事件组~\ref{def:complete-event-group}
  \item 概率的公理化定义~\ref{def:probability-axiom}
  \item 概率的基本性质~\ref{prop:probability-basic}
  \item 古典概型~\ref{def:classical-prob}
  \item 几何概型~\ref{def:geometric-prob}
  \item 条件概率~\ref{def:conditional-prob}
  \item 乘法公式~\ref{thm:multiplication}
  \item 全概率公式与贝叶斯公式~\ref{thm:total-probability}
  \item 事件的独立性~\ref{def:independence}
  \item 伯努利概型~\ref{def:bernoulli}
\end{introduction}

\section{基本概念}

\begin{definition}[随机试验] \label{def:random-experiment}
如果一个试验满足以下三个条件：
\begin{enumerate}
    \item 试验可以在相同的条件下重复进行；
    \item 试验所有可能结果明确可知，且不止一个；
    \item 每一次试验会出现哪一个结果，事先并不能确定，
\end{enumerate}
则称此试验为\textbf{随机试验}。随机试验简称试验，用字母 \(E\) 或 \(E_1, E_2, \dots\) 表示。
\end{definition}

我们是通过研究随机试验来研究随机现象的。

\begin{note}
在不少情况下，不能确切知道某一随机试验的全部可能结果，但可以知道它不超出某个范围。这时，也可以用这个范围来作为该试验的全部可能结果的集合。例如，需要记录某个城市一天的交通事故数量，则试验结果将是非负整数 \(x\)。无法确定 \(x\) 的可能取值的确切范围，但可以把这个范围取为 \([0, +\infty)\)，它总能包含一切可能的试验结果，尽管明知某些结果，如 \(x > 10000\) 是不会出现的。甚至把这个范围取为 \((-\infty, +\infty)\) 也无妨。这里体现了一定的数学抽象，它可以带来很大的方便。
\end{note}

\begin{definition}[样本空间] \label{def:sample-space}
随机试验的每一个可能结果称为\textbf{样本点}，记为 \(\omega\)。样本点的全体组成的集合称为\textbf{样本空间}（或基本事件空间），记为 \(\Omega\)（或 \(S\)），即 \(\Omega = \{\omega\}\)。
\end{definition}

由一个样本点构成的事件称为\textbf{基本事件}。

\begin{definition}[随机事件] \label{def:random-event}
在一次试验中可能出现，也可能不出现的结果称为\textbf{随机事件}，简称事件，用大写字母 \(A, B, C\) 等表示。随机事件是样本空间 \(\Omega\) 的子集。在每次试验中，当且仅当这一子集中的一个样本点出现时，称这一事件发生。
\end{definition}

\begin{enumerate}
    \item 对于仅由一个基本结果构成的单点集，称为\textbf{基本事件}；
    \item 将每次试验中一定发生的事件称为\textbf{必然事件}，记为 \(\Omega\)；
    \item 每次试验中一定不发生的事件称为\textbf{不可能事件}，记为 \(\varnothing\)。
\end{enumerate}

\begin{note}
概念之间的联系：随机试验产生基本结果，基本结果 \(\in\) 随机事件 \(\subset\) 样本空间。随机事件在一次试验中是否发生虽然不能事先确定，但是在大量重复试验的情况下，它的发生呈现出一定的规律性，这门课程正是要研究这种规律性。
\end{note}

\section{事件的关系与运算}

事件是一个集合，因此事件之间的关系和运算与集合论中集合之间的关系和运算完全类似，可以利用集合论的知识来理解事件的关系与运算。

\subsection{事件的关系}

\begin{enumerate}
    \item \textbf{包含}——如果事件 \(A\) 发生必导致事件 \(B\) 发生，则称事件 \(B\) 包含事件 \(A\)（或 \(A\) 被 \(B\) 包含），记为 \(A \subset B\)。也称 \(A\) 是 \(B\) 的子事件。
    \item \textbf{相等}——如果 \(A \subset B\) 且 \(B \subset A\)，则称事件 \(A\) 与 \(B\) 相等，记为 \(A = B\)。\(A\) 与 \(B\) 相等，事实上也就是说，\(A\) 与 \(B\) 由一些完全相同的试验结果构成，它不过是同一事件表面上看起来不同的两个说法而已。
    \item \textbf{相容}——若 \(AB \neq \varnothing\)，则称事件 \(A\) 和 \(B\) 相容。
    \item \textbf{互斥（互不相容）}——若 \(AB = \varnothing\)，则称事件 \(A\) 与 \(B\) 互不相容，也叫互斥。如果一些事件中任意两个事件都互斥，则称这些事件是两两互斥的。任意一对基本事件都是互斥的。
    \item \textbf{对立}——事件 \(A, B\) 不能同时发生，且至少一个发生，称事件 \(A, B\) 互为对立事件，或互为逆事件，即 \(A \cap B = \varnothing\)，\(A \cup B = \Omega\)。
\end{enumerate}

\begin{definition}[完备事件组] \label{def:complete-event-group}
如果 \(\bigcup_{i=1}^{n} A_i = \Omega\)（或 \(\bigcup_{i=1}^{\infty} A_i = \Omega\)），且 \(A_i A_j = \varnothing\)（对一切 \(i \neq j\)；\(i, j = 1, 2, \dots, n(\dots)\)），称有限个（或可列个）事件 \(A_1, A_2, \dots, A_n(\dots)\) 构成一个\textbf{完备事件组}。
\end{definition}

完备事件组在概率论中具有重要作用，全概率公式和贝叶斯公式都需要完备事件组作为前提条件。

\begin{note}
互斥与对立的区别：互斥只要求 \(AB = \varnothing\)，而对立还要求 \(A \cup B = \Omega\)。因此对立一定是互斥的，但互斥不一定是对立的。
\end{note}

\subsection{事件的运算}

\begin{enumerate}
    \item \textbf{积事件（交）}——称"事件 \(A\) 与 \(B\) 同时发生"的事件为事件 \(A\) 与 \(B\) 的\textbf{积事件}（或交事件），记为 \(A \cap B\) 或 \(AB\)。称"有限个（或可列个）事件 \(A_1, A_2, \dots, A_n(\dots)\) 同时发生"的事件为这些事件的积事件，记为 \(\bigcap_{i=1}^{n} A_i\)（或 \(\bigcap_{i=1}^{\infty} A_i\)）。
    \item \textbf{和事件（并）}——称"事件 \(A\) 与 \(B\) 至少有一个发生"的事件为事件 \(A\) 与 \(B\) 的\textbf{和事件}（或并事件），记为 \(A \cup B\)。称"有限个（或可列个）事件 \(A_1, A_2, \dots, A_n(\dots)\) 至少有一个发生"的事件为这些事件的和事件，记为 \(\bigcup_{i=1}^{n} A_i\)（或 \(\bigcup_{i=1}^{\infty} A_i\)）。
    \item \textbf{差事件}——称"事件 \(A\) 发生而事件 \(B\) 不发生"的事件为事件 \(A\) 与 \(B\) 的\textbf{差事件}，记为 \(A - B\)。
    \item \textbf{逆事件（对立事件）}——称"事件 \(A\) 不发生"的事件为事件 \(A\) 的\textbf{逆事件}或\textbf{对立事件}，记为 \(\overline{A}\)。
\end{enumerate}

由定义易知
\[
A - B = A - AB = A\overline{B}, \quad B = \overline{A} \Leftrightarrow AB = \varnothing \text{ 且 } A \cup B = \Omega.
\]

\subsection{事件的运算律}

\begin{enumerate}
    \item \textbf{吸收律}——若 \(A \subset B\)，则 \(A \cup B = B\)，\(A \cap B = A\)。
    \item \textbf{交换律}——\(A \cup B = B \cup A\)，\(A \cap B = B \cap A\)。
    \item \textbf{结合律}——\((A \cup B) \cup C = A \cup (B \cup C)\)，\((A \cap B) \cap C = A \cap (B \cap C)\)。
    \item \textbf{分配律}——\(A \cap (B \cup C) = (A \cap B) \cup (A \cap C)\)，\(A \cup (B \cap C) = (A \cup B) \cap (A \cup C)\)。特别地，
    \[
    A \cap (B - C) = (A \cap B) - (A \cap C).
    \]
    \item \textbf{对偶律（德·摩根律）}——\(\overline{A \cup B} = \overline{A} \cap \overline{B}\)，\(\overline{A \cap B} = \overline{A} \cup \overline{B}\)。一般地，对于 \(n \ge 1\)，
    \[
    \overline{\bigcup_{i=1}^{n} A_i} = \bigcap_{i=1}^{n} \overline{A_i}; \quad \overline{\bigcap_{i=1}^{n} A_i} = \bigcup_{i=1}^{n} \overline{A_i}.
    \]
\end{enumerate}

\begin{note}
德摩根律记忆技巧：长线变短线，开口换方向。事件运算顺序约定为先进行逆运算，然后进行交运算，最后进行并或差运算。事件的关系、运算与集合的关系、运算相当，且具有相同的运算法则，所以我们可以对比着理解记忆，并要学会用集合关系去考虑事件关系。
\end{note}

\begin{exercise}\label{exer:1-1}
已知 \(A\)、\(B\)、\(C\) 为三个随机事件，则 \(A\)、\(B\)、\(C\) 不都发生的事件为（\quad ）。\\
A. $\overline{A}\overline{B}\overline{C}$\quad B. $\overline{ABC}$\quad C. $A \cup B \cup C$\quad D. $ABC$
\end{exercise}

\begin{solution}
"\(A\)、\(B\)、\(C\) 不都发生"的对立事件为"\(A\)、\(B\)、\(C\) 都发生"，即 \(ABC\)，故"\(A\)、\(B\)、\(C\) 不都发生"为 \(\overline{ABC}\)。答案选 B。
\end{solution}

\begin{exercise}\label{exer:1-2}
设有随机事件 \(A_i\ (i=1,2,\dots,n)\)，\(A = \bigcap_{i=1}^{n} A_i\)，则 \(A\) 的对立事件 \(\overline{A} =\ \underline{\hspace{3cm}}\)。
\end{exercise}

\begin{solution}
根据德摩根律，
\[
\overline{A} = \overline{\bigcap_{i=1}^{n} A_i} = \overline{A_1 \cap A_2 \cap \dots \cap A_n} = \overline{A_1} \cup \overline{A_2} \cup \dots \cup \overline{A_n} = \bigcup_{i=1}^{n} \overline{A_i}.
\]
\end{solution}

\begin{exercise}\label{exer:1-28}
设 \(A,B,C\) 是任意三个事件，则下列选项中正确的是（\quad）。\par
A. 若 \(A\cup C=B\cup C\)，则 \(A=B\)\quad
B. 若 \(A-C=B-C\)，则 \(A=B\)\quad
C. 若 \(AC=BC\)，则 \(A=B\)\quad
D. 若 \(AB=\varnothing\) 且 \(\overline{A}\,\overline{B}=\varnothing\)，则 \(\overline{A}=B\)
\end{exercise}

\begin{solution}
应选 D。

\textbf{方法一}（直接法）：\(\overline{A}\,\overline{B}=\overline{A\cup B}=\varnothing\)，故 \(A\cup B=\Omega\)。又 \(AB=\varnothing\)，因此 \(A,B\) 互为对立事件，即 \(\overline{A}=B\)。

\textbf{方法二}（排除法）：取 \(C=\Omega\)，则 \(A\cup C=B\cup C=\Omega\) 恒成立，但 \(A\neq B\) 仍可成立，故 (A) 错。取 \(C=\varnothing\)，则 \(AC=BC=\varnothing\) 恒成立，故 (C) 错。(B) 类似取 \(C=\Omega\) 即可排除。
\end{solution}

\begin{note}
本题反映了事件运算与数的运算的不同之处。事件的并、交、差运算不满足消去律。
\end{note}

\begin{exercise}\label{exer:1-29}
用事件 \(A,B,C\) 的运算关系表示下列事件：\par
(1) \(A,B,C\) 都不发生；\par
(2) \(A,B,C\) 不都发生；\par
(3) \(A,B,C\) 不多于一个发生。
\end{exercise}

\begin{solution}
(1) \(A,B,C\) 都不发生 = 三者均不发生：
\[\overline{A}\,\overline{B}\,\overline{C} = \overline{A\cup B\cup C}.\]

(2) \(A,B,C\) 不都发生 = "都发生"的逆事件：
\[\overline{ABC} = \overline{A}\cup\overline{B}\cup\overline{C}.\]

(3) \(A,B,C\) 不多于一个发生 = 至多一个发生 = 至少有两个不发生：
\[\overline{A}\,\overline{B}\cup\overline{A}\,\overline{C}\cup\overline{B}\,\overline{C} = A\overline{B}\,\overline{C}\cup\overline{A}\,B\,\overline{C}\cup\overline{A}\,\overline{B}\,C\cup\overline{A}\,\overline{B}\,\overline{C}.\]
\end{solution}

\begin{exercise}\label{exer:1-30}
设随机变量 \(X,Y\)，事件 \(A=\{X\ge c\}\)，事件 \(B=\{Y\ge c\}\)。用 \(A,B\) 表示下列事件：\par
\(D=\{\min(X,Y)<c\}\)；\quad \(E=\{\max(X,Y)\ge c\}\)；\quad \(F=\{\min(X,Y)\ge c\}\)；\quad \(G=\{\max(X,Y)<c\}\)。
\end{exercise}

\begin{solution}
\begin{itemize}
    \item \(D = \overline{A}\cup\overline{B}\)（\(X,Y\) 中至少有一个 \(<c\)）；
    \item \(E = A\cup B\)（\(X,Y\) 中至少有一个 \(\ge c\)）；
    \item \(F = A\cap B\)（\(X,Y\) 均需 \(\ge c\)）；
    \item \(G = \overline{A}\cap\overline{B} = \overline{A\cup B}\)（\(X,Y\) 均需 \(<c\)）。
\end{itemize}
\end{solution}

\begin{exercise}\label{exer:1-31}
设 \(A\) 和 \(B\) 是任意两个事件，则下列两个命题（\quad）。\par
(1) 若 \(P(A)=P(B)\)，则 \(A=B\)；\par
(2) 若 \(P(AB)=0\)，则 \(AB=\varnothing\)。
\par
A. (1) 正确，(2) 不正确\quad
B. (1) 不正确，(2) 正确\quad
C. (1)、(2) 均正确\quad
D. (1)、(2) 均不正确
\end{exercise}

\begin{solution}
应选 D。

反例：向 \([0,2]\) 上随机投点。令 \(A=\{\text{点落入 }[0,1]\}\)，\(B=\{\text{点落入 }[1,2]\}\)。则 \(P(A)=P(B)=0.5\)，但 \(A\neq B\)；\(P(AB)=P(\{1\})=0\)，但 \(AB=\{1\}\neq\varnothing\)。故两个命题均不正确。
\end{solution}

\begin{note}
\(P(A)=0\) 不能推出 \(A=\varnothing\)，\(P(A)=1\) 不能推出 \(A=\Omega\)。概率只反映可能性大小，不反映事件的集合构成。
\end{note}

\begin{exercise}\label{exer:1-32}
已知事件 \(A,B\) 满足 \(P(AB)=P(\overline{A}\,\overline{B})\)，记 \(P(A)=p\)，则 \(P(B)=\underline{\hspace{2cm}}\)。
\end{exercise}

\begin{solution}
应填 \(1-p\)。

由德摩根律和加法公式：
\[P(\overline{A}\,\overline{B}) = P(\overline{A\cup B}) = 1-P(A\cup B) = 1-P(A)-P(B)+P(AB),\]
代入 \(P(AB)=P(\overline{A}\,\overline{B})\)，得
\[1-P(A)-P(B)=0 \implies P(B)=1-P(A)=1-p.\]
\end{solution}

\begin{exercise}\label{exer:1-33}
已知 \(P(A)=0.7\)，\(P(A-B)=0.3\)，则 \(P(\overline{AB})=\underline{\hspace{2cm}}\)。
\end{exercise}

\begin{solution}
应填 \(0.6\)。

由减法公式 \(P(A-B)=P(A)-P(AB)\)，得 \(P(AB)=0.7-0.3=0.4\)，故
\[P(\overline{AB}) = 1-P(AB) = 0.6.\]
\end{solution}

\begin{exercise}\label{exer:1-34}
设 \(A,B\) 为两个随机事件，且 \(P(A)=0.4\)，\(P(A\cup B)=0.7\)。若 \(A,B\) 互不相容，则 \(P(B)=\underline{\hspace{2cm}}\)；若 \(A,B\) 相互独立，则 \(P(B)=\underline{\hspace{2cm}}\)；若 \(A\) 发生则 \(B\) 必发生，则 \(P(B)=\underline{\hspace{2cm}}\)。
\end{exercise}

\begin{solution}
应填 \(0.3\)；\(0.5\)；\(0.7\)。

由加法公式 \(P(A\cup B)=P(A)+P(B)-P(AB)\)：
\begin{itemize}
    \item 互不相容时 \(P(AB)=0\)：\(P(B)=0.7-0.4=0.3\)；
    \item 相互独立时 \(P(AB)=P(A)\,P(B)\)：\(0.7=0.4+P(B)-0.4P(B)\)，\(P(B)=\dfrac{0.3}{0.6}=0.5\)；
    \item \(A\) 发生 \(B\) 必发生即 \(A\subset B\)：\(P(B)=P(A\cup B)=0.7\)。
\end{itemize}
\end{solution}

\section{概率的定义}

\subsection{描述性定义}

通常将随机事件 \(A\) 发生的可能性大小的度量（非负值）称为事件 \(A\) 发生的\textbf{概率}，记为 \(P(A)\)。

\subsection{统计性定义}

在相同条件下做重复试验，事件 \(A\) 出现的次数 \(k\) 和总的试验次数 \(n\) 之比 \(\frac{k}{n}\) 称为事件 \(A\) 在这 \(n\) 次试验中出现的\textbf{频率}。当试验次数 \(n\) 充分大时，频率将"稳定"于某常数 \(p\)。\(n\) 越大，频率偏离这个常数 \(p\) 的可能性越小。这个常数 \(p\) 就称为事件 \(A\) 的\textbf{概率}。

\begin{note}
概率的统计性定义实质上是说，用频率 \(\frac{k}{n}\) 作为事件 \(A\) 的概率 \(P(A)\) 的估计。其直观理解为某事件出现的可能性大小，可由其在多次重复试验中出现的频率去刻画。从上述可以看出，频率只是概率的估计，而非概率本身。概率的统计性定义的重要性主要基于以下两点：
\begin{enumerate}
    \item 它提供了估计概率的方法。比如在一批产品中抽取样品，来估计该批产品的合格率（合格率是客观的数据，抽取样品计算出来的合格率只是一种估计）。
    \item 它提供了一种检验某结论是否正确的准则。比如，你说某批产品的合格率是 95\%，我们做试验，抽取样品进行计算，得出的结果是合格率为 20\%，远远低于你所说的 95\%，于是毫不犹豫地拒绝你的结论。
\end{enumerate}
\end{note}

\subsection*{公理化定义}

\begin{definition}[概率的公理化定义] \label{def:probability-axiom}
设随机试验的样本空间为 \(\Omega\)，如果对每一个事件 \(A\) 都有一个确定的实数 \(P(A)\)，且事件函数 \(P(\cdot)\) 满足：
\begin{enumerate}
    \item \textbf{非负性}：\(P(A) \ge 0\)；
    \item \textbf{规范性}：\(P(\Omega) = 1\)；
    \item \textbf{可列可加性}：对任意可列个两两互不相容事件 \(A_1, A_2, \dots, A_n, \dots\)（即 \(A_i A_j = \varnothing\)，\(i \neq j\)；\(i, j = 1, 2, \dots\)），有
    \[
    P\!\left(\bigcup_{i=1}^{\infty} A_i\right) = \sum_{i=1}^{\infty} P(A_i),
    \]
\end{enumerate}
则称 \(P(\cdot)\) 为\textbf{概率}，\(P(A)\) 为事件 \(A\) 的\textbf{概率}。
\end{definition}

\begin{remark}
(1) 数学上所说的"公理"，就是一些不加证明而承认的前提，上述公理化定义只是界定了概率这个概念所必须满足的一般性质，它不解决具体场合下的概率计算。\newline
(2) 概率 \(P(\cdot)\) 是事件的函数。\newline
(3) 虽然公理化定义不解决具体场合下的概率计算，但是我们却经常用它来判断事件函数 \(P(\cdot)\) 是否是概率，这种题型在考研试题中也会经常遇到。
\end{remark}

由概率的公理化定义，可以推导出以下基本性质：

\begin{property}\label{prop:probability-basic}
概率的基本性质：
\begin{enumerate}
    \item \(P(\varnothing) = 0\)；
    \item \textbf{有限可加性}：若 \(A_1, A_2, \dots, A_n\) 两两互不相容，则
    \[
    P\!\left(\bigcup_{i=1}^{n} A_i\right) = \sum_{i=1}^{n} P(A_i);
    \]
    \item \textbf{减法公式}：若 \(B \subset A\)，则 \(P(A - B) = P(A) - P(B)\)。更一般地，
    \[
    P(A - B) = P(A) - P(AB);
    \]
    \item \textbf{单调性}：若 \(B \subset A\)，则 \(P(B) \le P(A)\)；
    \item \textbf{有界性}：对任意事件 \(A\)，\(0 \le P(A) \le 1\)；
    \item \textbf{求逆公式}：\(P(\overline{A}) = 1 - P(A)\)；
    \item \textbf{加法公式}：\(P(A \cup B) = P(A) + P(B) - P(AB)\)。当 \(A\)、\(B\) 互斥时，\(P(A \cup B) = P(A) + P(B)\)。
    \item 设 \(A_1, A_2, A_3\) 为任意三个事件，则有
    \[
    P(A_1 \cup A_2 \cup A_3) = P(A_1) + P(A_2) + P(A_3) - P(A_1A_2) - P(A_1A_3) - P(A_2A_3) + P(A_1A_2A_3).
    \]
    \item 若 \(A_1, A_2, \dots, A_n\) 是两两互不相容的事件，则有
    \[
    P(A_1 \cup A_2 \cup \dots \cup A_n) = P(A_1) + P(A_2) + \dots + P(A_n).
    \]
\end{enumerate}
\end{property}
\begin{note}
概率为 0 并不意味着为空集。即 \(P(A) = 0 \nRightarrow A = \varnothing\)，而 \(A = \varnothing \Rightarrow P(A) = 0\)。同样的，$P(A) = 0$不能断言$A = \varnothing$，而$P(A) = 1$不能断言$A = \Omega$.
\end{note}

\begin{exercise}\label{exer:1-3}
设 \(A, B\) 为两个随机事件，\(P(A) = 0.4\)，\(P(A \cup B) = 0.7\)。\par
(1) 求 \(P(B - A)\)。\par
(2) 若事件 \(A\) 和事件 \(B\) 互不相容，求 \(P(B)\)。
\end{exercise}

\begin{solution}
(1) 由加法公式
\[
P(A \cup B) = P(A) + P(B) - P(AB) = P(A) + P(B - A),
\]
故
\[
P(B - A) = P(A \cup B) - P(A) = 0.7 - 0.4 = 0.3.
\]

(2) 若 \(A\) 与 \(B\) 互不相容，则 \(P(AB) = 0\)，此时加法公式简化为
\[
P(A \cup B) = P(A) + P(B),
\]
故
\[
P(B) = P(A \cup B) - P(A) = 0.7 - 0.4 = 0.3.
\]
\end{solution}


\begin{exercise}\label{exer:1-27}
设 \(A\) 和 \(B\) 是任意两个概率不为 0 的互不相容事件，则下列结论中肯定正确的是（\quad）。\par
A. \(\overline{A}\) 与 \(\overline{B}\) 不相容\quad
B. \(\overline{A}\) 与 \(\overline{B}\) 相容\quad
C. \(P(AB) = P(A)P(B)\)\quad
D. \(P(A-B) = P(A)\)
\end{exercise}

\begin{solution}
应选 D。

由 \(A\) 与 \(B\) 互不相容，\(P(AB) = 0\)，故
\[P(A-B) = P(A) - P(AB) = P(A).\]

当 \(A\) 与 \(B\) 互为对立事件时，\(\overline{A}\) 与 \(\overline{B}\) 不相容；反之则相容。故选项 A、B 不一定正确。选项 C：\(P(AB) = 0 \neq P(A)P(B)\)（因为 \(P(A)>0, P(B)>0\)），故错误。
\end{solution}
\begin{exercise}\label{exer:1-35}
考虑一元二次方程 \(x^2+Bx+C=0\)，其中 \(B,C\) 分别是将一枚骰子接连抛两次先后出现的点数。则该方程有实根的概率为（\quad）。\par
A. \(\dfrac{19}{36}\)\quad
B. \(\dfrac{17}{36}\)\quad
C. \(\dfrac{15}{36}\)\quad
D. \(\dfrac{13}{36}\)
\end{exercise}

\begin{solution}
应选 A。

方程有实根的条件为 \(B^2-4C\ge 0\)。总样本数 \(6^2=36\)。列表统计满足条件的样本点数：
\[
\begin{array}{c|cccccc}
B\setminus C & 1 & 2 & 3 & 4 & 5 & 6 \\
\midrule
1 & - & - & - & - & - & - \\
2 & 0 & - & - & - & - & - \\
3 & + & + & - & - & - & - \\
4 & + & + & + & 0 & - & - \\
5 & + & + & + & + & + & + \\
6 & + & + & + & + & + & + \\
\end{array}
\]
``+''和``0''对应的样本点共 \(1+2+4+6+6=19\) 个，故 \(P=\dfrac{19}{36}\)。
\end{solution}

\begin{exercise}\label{exer:1-36}
甲口袋有 5 个白球、3 个黑球，乙口袋有 4 个白球、6 个黑球。从两个口袋中各任取一个球，则取到的两个球颜色相同的概率为（\quad）。\par
A. \(\dfrac{19}{40}\)\quad
B. \(\dfrac{21}{40}\)\quad
C. \(\dfrac{1}{4}\)\quad
D. \(\dfrac{9}{40}\)
\end{exercise}

\begin{solution}
应选 A。

两个球颜色相同分两种情况：都白或都黑。总取法 \(8\times10=80\) 种，有利取法 \(5\times4+3\times6=38\) 种，故
\[P = \frac{5\times4+3\times6}{8\times10} = \frac{38}{80} = \frac{19}{40}.\]
\end{solution}

\begin{exercise}\label{exer:1-45}
已知 \(P(B)=0.5\)，\(P(A\cup B)=0.7\)，则 \(P(A\overline{B})=\underline{\hspace{2cm}}\)，\(P(\overline{A}\overline{B})=\underline{\hspace{2cm}}\)。
\end{exercise}

\begin{solution}
应填 \(0.2\)；\(0.3\)。

由加法公式：
\[
P(A\overline{B}) = P(A) - P(AB) = P(A\cup B) - P(B) = 0.7 - 0.5 = 0.2,
\]
\[
P(\overline{A}\overline{B}) = P(\overline{A\cup B}) = 1 - P(A\cup B) = 1 - 0.7 = 0.3.
\]
\end{solution}

\begin{exercise}\label{exer:1-46}
对任意事件 \(A,B\)，下面结论正确的是（\quad）。\par
A. 若 \(P(AB)=0\)，则 \(\overline{A}\,\overline{B}=\Omega\)\quad
B. \(P(A-B)=P(A)-P(B)\)\quad
C. \(P(A\overline{B})=P(A)-P(B)+P(\overline{A}B)\)\quad
D. \(P(A\cup B)=P(A)+P(B)-P(A)P(B)\)
\end{exercise}

\begin{solution}
应选 C。

A 错误：\(P(AB)=0\) 不能推出 \(AB=\varnothing\)，从而不能推出 \(\overline{A}\,\overline{B}=\Omega\)。B 错误：\(P(A-B)=P(A)-P(AB)\)，一般不等于 \(P(A)-P(B)\)。D 错误：题目未说明 \(A,B\) 独立，故 \(P(AB)\neq P(A)P(B)\)。

C 正确：\(P(A\overline{B})=P(A)-P(AB)\)，\(P(\overline{A}B)=P(B)-P(AB)\)，两式消去 \(P(AB)\) 得
\[
P(A\overline{B}) = P(A) - P(B) + P(\overline{A}B).
\]
\end{solution}

\begin{exercise}\label{exer:1-47}
设 \(A\) 和 \(B\) 是任意两个随机事件，则（\quad）。\par
A. \(P(A+B)\le P(A\cup B)\le P(A)+P(B)\)\quad
B. \(P(A)+P(B)-1\le P(AB)\le P(A\cup B)\le P(A)+P(B)\)\quad
C. \(P(A)+P(B)-1\le P(A(A\cup B))\le P(AB)\le P(A\cup(A\cup B))\)\quad
D. \(1-P(A)-P(B)\le P(AB)\le P(A)+P(B)\le P(A\cup B)\)
\end{exercise}

\begin{solution}
应选 B。

A 错误：\(P(A+B)=P(A\cup B)\)（概率中 "+" 即并）。C 错误：\(P(A(A\cup B))=P(A)\ge P(AB)\)。D 错误：\(P(A)+P(B)\ge P(A\cup B)\)，故 \(P(A)+P(B)\le P(A\cup B)\) 不成立。

B 正确：由 \(P(AB)=P(A)+P(B)-P(A\cup B)\)：
\begin{itemize}
    \item 左边：\(P(A\cup B)\le 1\)，故 \(P(AB)\ge P(A)+P(B)-1\)；
    \item 中间：\(AB\subset A\cup B\)，故 \(P(AB)\le P(A\cup B)\)；
    \item 右边：\(P(A\cup B)\le P(A)+P(B)\)（加法公式的直接推论）。
\end{itemize}
\end{solution}

\begin{exercise}\label{exer:1-48}
设 \(A\) 与 \(B\) 是两个随机事件，且 \(P(AB)=0\)，则（\quad）。\par
A. \(P(A-B)=P(A)\)\quad
B. \(A\) 与 \(B\) 相互独立\quad
C. \(P(A)=0\) 或 \(P(B)=0\)\quad
D. \(A\) 与 \(B\) 互不相容
\end{exercise}

\begin{solution}
应选 A。

\(P(A-B)=P(A)-P(AB)=P(A)\)，故 A 正确。B 与条件无必然联系；C 错误（例如连续型随机变量取特定值的概率为零，但事件本身可能发生）；D 错误（\(P(AB)=0\) 是互不相容的必要条件而非充分条件）。
\end{solution}


\begin{exercise}[提高]\label{exer:1-49}
设事件 \(A,B\) 相互独立，事件 \(B,C\) 互不相容，事件 \(A\) 与 \(C\) 不能同时发生，且 \(P(A)=P(B)=0.5\)，\(P(C)=0.2\)。则事件 \(A,B\) 和 \(C\) 中仅 \(C\) 发生或仅 \(C\) 不发生的概率为 \underline{\hspace{2cm}}。
\end{exercise}

\begin{solution}
应填 \(0.45\)。

"仅 \(C\) 发生"即 \(C\overline{A}\,\overline{B}\)。由于 \(AC=\varnothing\) 且 \(BC=\varnothing\)，故 \(C\subset\overline{A}\) 且 \(C\subset\overline{B}\)，即 \(C\subset\overline{A}\,\overline{B}\)，从而 \(C\overline{A}\,\overline{B}=C\)，\(P(C)=0.2\)。

"仅 \(C\) 不发生"即 \(\overline{C}AB\)。由独立性 \(P(AB)=P(A)P(B)=0.25\)，且 \(AB\subset\overline{C}\)（因为 \(AC=\varnothing\)），故 \(\overline{C}AB=AB\)，\(P(\overline{C}AB)=0.25\)。

又 \(C\) 与 \(AB\) 互斥（\(C\subset\overline{A}\)，\(AB\subset A\)），故
\[
P(C\overline{A}\,\overline{B}\cup\overline{C}AB) = P(C)+P(AB) = 0.2+0.25 = 0.45.
\]
\end{solution}

\begin{exercise}[提高]\label{exer:1-50}
设 \(A,B\) 为任意两个随机事件，则（\quad）。\par
A. \(P(AB)\le P(A)P(B)\)\quad
B. \(P(AB)\ge P(A)P(B)\)\quad
C. \(P(AB)\le \dfrac{P(A)+P(B)}{2}\)\quad
D. \(P(AB)\ge \dfrac{P(A)+P(B)}{2}\)
\end{exercise}

\begin{solution}
应选 C。

由 \(AB\subset A\) 和 \(AB\subset B\)，得 \(P(AB)\le P(A)\) 且 \(P(AB)\le P(B)\)，因此
\[
2P(AB)\le P(A)+P(B) \implies P(AB)\le\frac{P(A)+P(B)}{2}.
\]

A、B 的反例：当 \(A\subset B\) 时 \(P(AB)=P(A)\)，与 \(P(A)P(B)\) 无确定大小关系；当 \(AB=\varnothing\) 时 \(P(AB)=0\)。D 的反例同 A。
\end{solution}


\section{古典概型}

\begin{definition}[古典概型] \label{def:classical-prob}
如果随机试验的样本空间满足：
\begin{enumerate}
    \item 只有有限个样本点（基本事件）；
    \item 每个样本点（基本事件）发生的可能性都一样，
\end{enumerate}
则称随机试验的概率模型为\textbf{古典概型}（等可能概型）。
\end{definition}

如果古典概型的基本事件总数为 \(n\)，事件 \(A\) 包含 \(k\) 个基本事件，则 \(A\) 的概率为
\[
P(A) = \frac{k}{n} = \frac{\text{事件} \ A \ \text{所含基本事件的个数}}{\text{基本事件总数}}.
\]

由上式计算得出的概率称为 \(A\) 的\textbf{古典概率}。

\begin{note}
计算古典概率的关键是基本事件和样本空间的选定以及基本事件数的计算：
\begin{enumerate}
    \item \textbf{列举法}（直接数数法）：基本事件数不多时常用。
    \item \textbf{集合对应法}：
    \begin{enumerate}
        \item \textbf{加法原理}——完成一件事有 \(n\) 类办法，第一类办法中有 \(m_1\) 种方法，第二类办法中有 \(m_2\) 种方法，\(\dots\)，第 \(n\) 类办法中有 \(m_n\) 种方法，则完成此事共有 \(\sum_{i=1}^{n} m_i\) 种方法。
        \item \textbf{乘法原理}——完成一件事有 \(n\) 个步骤，第一步有 \(m_1\) 种方法，第二步有 \(m_2\) 种方法，\(\dots\)，第 \(n\) 步有 \(m_n\) 种方法，则完成此事共有 \(\prod_{i=1}^{n} m_i\) 种方法。
        \item \textbf{排列}——从 \(n\) 个不同的元素中取出 \(m(\le n)\) 个元素，并按照一定顺序排成一列，叫作排列。所有排列的个数叫作排列数，记作
        \[
        \mathrm{P}_n^m = n(n-1)(n-2)\cdots(n-m+1) = \frac{n!}{(n-m)!}.
        \]
        当 \(m=n\) 时，\(\mathrm{P}_n^n = \frac{n!}{0!} = n!\)，叫作全排列。
        \item \textbf{组合}——从 \(n\) 个不同的元素中取出 \(m(\le n)\) 个元素，并成一组，叫作组合。所有组合的个数叫作组合数，记作
        \[
        \mathrm{C}_n^m = \frac{\mathrm{P}_n^m}{m!} = \frac{n!}{m!(n-m)!}, \quad \mathrm{C}_n^m = \mathrm{C}_n^{n-m}, \quad 0!=1, \quad \mathrm{C}_n^0=1.
        \]
    \end{enumerate}
    \item \textbf{逆数法}：先求 \(\overline{A}\) 中的基本事件数 \(n_{\overline{A}}\)，用基本事件总数 \(n\) 减去 \(n_{\overline{A}}\) 便得 \(A\) 中的基本事件数，这种方法常用于计算含有"至少"字样的事件的概率。
\end{enumerate}
\end{note}

\subsection*{随机分配问题}

随机分配也叫随机占位，突出一个"放"字，即将 \(n\) 个可辨质点随机地分配到 \(N\) 个盒子中，区分每盒可以容纳任意多个质点（ 不同分法的总数\(N^n\)）和最多可以容纳一个质点（不同分法的总数\(\mathrm{P}_N^n = N(N-1)\cdots(N-n+1)\)）。

\begin{note}
每盒容纳任意多个质点：每个质点均可放到 \(N\) 个盒子中的任何一个，即有 \(N\) 种放法，于是 \(n\) 个可辨质点放到 \(N\) 个盒子中共有 \(N^n\) 种不同放法。\newline
每盒容纳至多一个质点：质点可辨，且一个盒子至多容纳一个质点，故 \(n\) 个质点放到 \(N(N\ge n)\) 个盒子中的所有不同放法即从 \(N\) 个元素中选取 \(n\) 个元素的排列数 \(\mathrm{P}_N^n\)。
\end{note}

\begin{exercise}\label{exer:1-4}
将 \(n\) 个球随机放入 \(N(n\le N)\) 个盒子中，每个盒子可以放任意多个球。求下列事件的概率：\par
(1) \(A=\{\text{某指定} \ n \ \text{个盒子各有一球}\}\)；\par
(2) \(B=\{\text{恰有} \ n \ \text{个盒子各有一球}\}\)；\par
(3) \(C=\{\text{指定} \ k(k\le n) \ \text{个盒子各有一球}\}\)。
\end{exercise}

\begin{solution}
设 \(n\) 个球、\(N\) 个盒子是可分辨的（例如编号），由于每个盒子可以放任意多个球，因此每个球都有 \(N\) 种不同的放置方法。将 \(n\) 个球随机放入 \(N\) 个盒子的一种放法作为基本事件，则基本事件总数为 \(N^n\)。

(1) 事件 \(A\) 所含的基本事件是 \(n\) 个不同球的一种排列，故
\[
n_A = n!, \quad P(A) = \frac{n!}{N^n}.
\]

(2) 事件 \(B\) 中的基本事件可以设想为先从 \(N\) 个盒子中选出 \(n\) 个（共有 \(\mathrm{C}_N^n\) 种不同选法），而后把 \(n\) 个球随机放入这 \(n\) 个盒子中，每盒一球（共有 \(n!\) 种不同放法），因此
\[
n_B = \mathrm{C}_N^n \cdot n!, \quad P(B) = \frac{\mathrm{C}_N^n \cdot n!}{N^n}.
\]

(3) 事件 \(C\) 的基本事件可以设想为先从 \(n\) 个球中选出 \(k\) 个球（有 \(\mathrm{C}_n^k\) 种不同选法），再将这 \(k\) 个球随机放入指定的 \(k\) 个盒子中，每盒一球（有 \(k!\) 种不同放法），最后将余下的 \(n-k\) 个球随机放入其余的 \(N-k\) 个盒子中，每个球都有 \(N-k\) 种放置方法，因此共有 \((N-k)^{n-k}\) 种不同放法，故
\[
n_C = \mathrm{C}_n^k \cdot k! \cdot (N-k)^{n-k}, \quad P(C) = \frac{\mathrm{C}_n^k k! (N-k)^{n-k}}{N^n} = \frac{n! (N-k)^{n-k}}{(n-k)! N^n}.
\]
\end{solution}

\begin{note}
许多问题的结构形式与分球入盒问题相同，都属于随机占位问题。例如生日问题（\(n\) 个人生日，相当于 \(n\) 个球随机放入 365 个盒子中，每盒可以放多个球）；住房分配问题（\(n\) 个人被分配到 \(N\) 个房间中去，每个房间可住多个人）；乘客下车问题（\(n\) 个乘客在 \(N\) 个车站下车的各种可能情况）；等等。对这些问题的求解都可以用"将 \(n\) 个球等可能地投放到 \(N\) 个盒子中"的思路来考虑。
\end{note}

\subsection*{简单随机抽样问题}

设 \(\Omega = \{\omega_1, \omega_2, \dots, \omega_N\}\) 含 \(N\) 个元素，称 \(\Omega\) 为总体。如果各元素被抽到的可能性相同，则总体 \(\Omega\) 的抽样称作\textbf{简单随机抽样}，突出一个"取"字。

简单随机抽样分为先后有放回、先后无放回及任取这三种不同的方式。

\begin{enumerate}
  \item 先后有放回取 \(n\) 次 ，抽取法总数 \(N^n\) 
  \item 先后无放回取 \(n\) 次 ，抽取法总数 \(\mathrm{P}_N^n = N(N-1)\cdots(N-n+1)\)
  \item 任取 \(n\) 个，抽取法总数 \(\mathrm{C}_N^n\)
\end{enumerate}

\begin{note}
先后有放回取 \(n\) 次：既考虑抽到何元素，又考虑各元素出现的顺序，每次从 \(\Omega\) 中随意抽取一个元素，并在抽下一元素前将其放回 \(\Omega\)，于是每次都有 \(N\) 个元素可被抽取，即有 \(N\) 种抽取方法，抽取 \(n\) 次，即 \(N^n\)。\newline
先后无放回取 \(n\) 次：既考虑抽到何元素，又考虑各元素出现的顺序，只是抽出的元素均不再放回 \(\Omega\)，于是每次抽取时都比上一次少了一个元素，抽 \(n(n\le N)\) 次，即 \(\mathrm{P}_N^n = N(N-1)\cdots(N-n+1)\)。\newline
任取 \(n\) 个：一次性取 \(n\) 个元素，相当于将 \(n\) 个元素无序且无放回地取走，共抽取法总数为 \(\mathrm{C}_N^n = \dfrac{\mathrm{P}_N^n}{n!}\)。
\end{note}

\begin{exercise}\label{exer:1-5}
袋中有 5 个球，3 个白球，2 个黑球。\par
(1) 先后有放回取 2 个球；\par
(2) 先后无放回取 2 个球；\par
(3) 任取 2 个球。\par
求取的 2 个球中至少 1 个是白球的概率。
\end{exercise}

\begin{solution}
用对立事件思想，计算"两球全黑"的种数，再用总数减去它。
\begin{enumerate}
    \item 先后有放回取 2 个球：
    总数 \(5^2 = 25\)，两球全黑 \(2^2 = 4\)，故
    \[
    P = \frac{5^2 - 2^2}{5^2} = \frac{21}{25}.
    \]
    \item 先后无放回取 2 个球：
    总数 \(5 \cdot 4 = 20\)，两球全黑 \(2 \cdot 1 = 2\)，故
    \[
    P = \frac{5 \cdot 4 - 2 \cdot 1}{5 \cdot 4} = \frac{18}{20} = \frac{9}{10}.
    \]
    \item 任取 2 个球：
    总数 \(\mathrm{C}_5^2 = 10\)，两球全黑 \(\mathrm{C}_2^2 = 1\)，故
    \[
    P = \frac{\mathrm{C}_5^2 - \mathrm{C}_2^2}{\mathrm{C}_5^2} = \frac{9}{10}.
    \]
\end{enumerate}
\end{solution}

\begin{note}
上述 (2) 与 (3) 概率相同，这是因为 \(5 \cdot 4 = \mathrm{C}_5^2 \cdot 2!\)，左边上下有序，右边上下无序，相当于把顺序"消掉"了，故"先后无放回取 \(k\) 个球"与"任取 \(k\) 个球"的概率相同。由于 (3) 较方便，因此计算 (2) 时可按 (3) 来计算。
\end{note}

\begin{exercise}\label{exer:1-6}
袋中有 100 个球，40 个白球，60 个黑球。\par
(1) 先后有放回取 20 个球，求取出 15 个白球、5 个黑球的概率；\par
(2) 先后无放回取 20 个球，求取出 15 个白球、5 个黑球的概率；\par
(3) 先后有放回取 20 个球，求第 20 次取到白球的概率；\par
(4) 先后无放回取 20 个球，求第 20 次取到白球的概率。
\end{exercise}

\begin{solution}
\begin{enumerate}
    \item \[
    p_1 = \frac{\mathrm{C}_{20}^{15} \cdot 40^{15} \cdot 60^{5}}{100^{20}}.
    \]
    \item 按照"任取 20 个球"来计算，即
    \[
    p_2 = \frac{\mathrm{C}_{40}^{15} \mathrm{C}_{60}^{5}}{\mathrm{C}_{100}^{20}}.
    \]
    \item 有放回取球，每次抽取的样本空间没有变化，故每次取到白球的概率始终为 \(\dfrac{40}{100} = \dfrac{2}{5}\)，即
    \[
    p_3 = \frac{40}{100} = \frac{2}{5}.
    \]
    \item 看作随机占位问题：设有 100 个盒子，每个盒子中放 1 个球，则要求第 20 个盒子中放入白球即可。故
    \[
    p_4 = \frac{\mathrm{C}_{40}^1 \cdot 99!}{100!} = \frac{40}{100} = \frac{2}{5}.
    \]
\end{enumerate}
\end{solution}

\begin{note}
\(p_4 = p_3\)。无放回抽取时，每次取到白球的概率也不会变，这可理解成依概率摸球（抓阄模型）。类似地，设有 100 个灰球，白的成分与黑的成分之比为 \(40:60\)，每次取到球中白的成分始终是 \(40\%\)，不受抽取顺序的影响。
\end{note}

\begin{exercise}\label{exer:1-26}
从一副扑克牌（52 张）中任取 3 张（不重复），则取出 3 张牌中至少有 2 张花色相同的概率为 \underline{\hspace{2cm}}。
\end{exercise}

\begin{solution}
应填 \(\dfrac{256}{425}\)。

利用对立事件：所求概率的对立事件为"3 张牌花色各不相同"。从 4 种花色中选 3 种，每种花色取 1 张：
\[P(\text{花色各不相同}) = \frac{\mathrm{C}_4^3 \times 13^3}{\mathrm{C}_{52}^3} = \frac{4 \times 2197}{22100} = \frac{8788}{22100}.\]
因此
\[P(\text{至少 2 张同花色}) = 1 - \frac{8788}{22100} = \frac{13312}{22100} = \frac{256}{425}.\]
\end{solution}


\begin{exercise}[提高]\label{exer:1-51}
从 1 到 9 的 9 个整数中有放回地随机取 3 次，每次取一个数，则取出的 3 个数之积能被 10 整除的概率为 \underline{\hspace{2cm}}。
\end{exercise}

\begin{solution}
应填 \(\dfrac{52}{243}\)。

积能被 10 整除当且仅当积同时含有因子 2 和因子 5。从 \(\{1,\dots,9\}\) 中，5 的倍数仅有 \(\{5\}\)，2 的倍数为 \(\{2,4,6,8\}\)。

设 \(A\) 为"至少取出一个 5 的倍数"，\(B\) 为"至少取出一个 2 的倍数"。由容斥原理：
\begin{align*}
P(AB) &= 1 - P(\overline{A}) - P(\overline{B}) + P(\overline{A}\,\overline{B}) \\
&= 1 - \frac{8^3}{9^3} - \frac{5^3}{9^3} + \frac{4^3}{9^3} \\
&= 1 - \frac{512}{729} - \frac{125}{729} + \frac{64}{729} = \frac{156}{729} = \frac{52}{243}.
\end{align*}
\end{solution}


\subsection{几何概型}

\begin{definition}[几何概型] \label{def:geometric-prob}
如果随机试验的样本空间满足：
\begin{enumerate}
    \item 样本空间（基本事件空间）\(\Omega\) 是一个可度量的有界区域；
    \item 每个样本点（基本事件）发生的可能性都一样，即样本点落入 \(\Omega\) 的某一可度量的子区域 \(S\) 的可能性大小与 \(S\) 的几何度量成正比，而与 \(S\) 的位置及形状无关，
\end{enumerate}
则称随机试验的概率模型为\textbf{几何概型}。
\end{definition}

如果 \(S_A\) 是样本空间 \(\Omega\) 的一个可度量的子区域，则事件 \(A = \{\text{样本点落入区域 } S_A\}\) 的概率为
\[
P(A) = \frac{S_A \text{ 的几何度量}}{\Omega \text{ 的几何度量}}.
\]

\begin{note}
古典概型与几何概型的区别：基本事件有限、等可能发生的随机试验为古典概型；基本事件无限且具有几何度量、等可能发生的随机试验为几何概型。
\end{note}

\begin{exercise}\label{exer:1-7}
在区间 \((0,1)\) 中随机地取两个数，则这两个数之差的绝对值小于 \(0.5\) 的概率为 \underline{\hspace{2cm}}。
\end{exercise}

\begin{solution}
设两个数分别为 \(x, y\)，则 \((x, y)\) 的取值范围为正方形区域 \(D = \{(x, y) \mid 0 < x < 1, 0 < y < 1\}\)。又 \(|x - y| < 0.5\)，则所在区域如图中阴影部分所示，于是所求概率为
\[
p = \frac{1^2 - 0.5^2}{1^2} = 0.75.
\]
\end{solution}

\begin{exercise}\label{exer:1-24}
长方形 \(ABCD\) 的长 \(AB=10\,\mathrm{cm}\)，宽 \(BC=5\,\mathrm{cm}\)，在长方形 \(ABCD\) 内任取一点 \(E\)，则 \(P\!\left(\angle AEB \ge \dfrac{\pi}{2}\right) = \underline{\hspace{2cm}}\)。
\end{exercise}

\begin{solution}
应填 \(\dfrac{\pi}{4}\)。

\(\angle AEB \ge 90^\circ\) 当且仅当点 \(E\) 位于以 \(AB\) 为直径的半圆内。该圆半径 \(r = \dfrac{AB}{2} = 5\,\mathrm{cm}\)，半圆面积 \(S_{\text{半圆}} = \dfrac{1}{2}\pi r^2 = \dfrac{25\pi}{2}\)。由于半圆半径恰好等于矩形宽 \(BC = 5\,\mathrm{cm}\)，半圆完全包含在矩形内，故
\[P = \frac{S_{\text{半圆}}}{S_{\text{矩形}}} = \frac{25\pi/2}{50} = \frac{\pi}{4}.\]
\end{solution}


\begin{exercise}\label{exer:1-37}
从 \([0,1]\) 中随机地取两个数，其积大于 \(\dfrac{1}{4}\)，其和小于 \(\dfrac{5}{4}\) 的概率为 \underline{\hspace{2cm}}。
\end{exercise}

\begin{solution}
应填 \(\dfrac{15}{32}-\dfrac{1}{2}\ln 2\)。

设两个数分别为 \(x,y\)，则 \((x,y)\) 在正方形 \(\Omega=[0,1]^2\) 内均匀分布。所求事件为
\[A = \left\{(x,y)\,\middle|\, x+y<\frac{5}{4},\ xy>\frac{1}{4}\right\}.\]
曲线 \(xy=\dfrac{1}{4}\) 与 \(x+y=\dfrac{5}{4}\) 交于 \(x=\dfrac{1}{4}\) 和 \(x=1\)。对 \(x\in\left[\dfrac{1}{4},1\right]\)，\(y\) 的范围从 \(\dfrac{1}{4x}\) 到 \(\dfrac{5}{4}-x\)：
\begin{align*}
P(A) &= \int_{1/4}^{1}\left(\frac{5}{4}-x-\frac{1}{4x}\right)\mathrm{d}x = \left[\frac{5}{4}x-\frac{1}{2}x^2-\frac{1}{4}\ln x\right]_{1/4}^{1} \\
&= \frac{3}{4}-\left(\frac{9}{32}+\frac{1}{2}\ln 4\right) = \frac{15}{32}-\frac{1}{2}\ln 2.
\end{align*}
\end{solution}

\begin{exercise}\label{exer:1-38}
随机地向半圆 \(0<y<\sqrt{2ax-x^2}\)（\(a>0\)）内投掷一点，求该点和原点连线与 \(x\) 轴的夹角 \(\theta\le\dfrac{\pi}{4}\) 的概率。
\end{exercise}

\begin{solution}
应填 \(\dfrac{1}{2}+\dfrac{1}{\pi}\)。

半圆方程 \(y=\sqrt{2ax-x^2}\) 即 \((x-a)^2+y^2=a^2\)（\(y>0\)），是以 \((a,0)\) 为圆心、\(a\) 为半径的上半圆。直线 \(y=x\) 与半圆交于 \((0,0)\) 和 \((a,a)\)。

\(\theta\le\pi/4\) 对应 \(y\le x\)。在半圆内，该区域由两部分组成：
\begin{itemize}
    \item 三角形区域 \((0,0),(a,0),(a,a)\)：面积为 \(\dfrac{a^2}{2}\)；
    \item 半圆从 \(x=a\) 到 \(x=2a\) 的部分：面积为 \(\dfrac{1}{4}\pi a^2\)。
\end{itemize}
半圆总面积 \(S_\Omega=\dfrac{1}{2}\pi a^2\)，故
\[P = \frac{a^2/2+\pi a^2/4}{\pi a^2/2} = \frac{1}{2}+\frac{1}{\pi}.\]
\end{solution}


\section{条件概率}

\subsection{条件概率的定义}

\begin{definition}[条件概率] \label{def:conditional-prob}
设 \(A, B\) 是两个事件，且 \(P(A) > 0\)，称
\[
P(B|A) = \frac{P(AB)}{P(A)}
\]
为在事件 \(A\) 发生的条件下事件 \(B\) 发生的\textbf{条件概率}。
\end{definition}

\begin{property}\label{prop:conditional-prob}
条件概率 \(P(\cdot|A)\) 满足概率的三条公理：
\begin{enumerate}
    \item \textbf{非负性}：对于每一事件 \(B\)，有 \(P(B|A) \ge 0\)；
    \item \textbf{规范性}：\(P(\Omega|A) = 1\)；
    \item \textbf{可列可加性}：设 \(B_1, B_2, \dots, B_n, \dots\) 是两两互斥的事件，则
    \[
    P\!\left(\bigcup_{i=1}^{\infty} B_i \Big| A\right) = \sum_{i=1}^{\infty} P(B_i|A).
    \]
\end{enumerate}
因此，概率的一切性质和公式对条件概率都适用。
\end{property}

\begin{note}
对于条件概率 \(P(B|A)\)，可以理解为样本空间由原来的 \(\Omega\) 缩小为 \(A\)，即
\[
P(B) = P(B|\Omega) \xrightarrow{\Omega \to A} P(B|A).
\]
从这个角度理解，样本空间的缩小使不确定性减小了，即当 \(A \subset B \subset \Omega\) 时，有 \(P(A|B) \ge P(A)\) 成立。条件概率就是在"已知某事件发生了"这一附加条件下所计算的概率。例如，
\[
P(\overline{B}|A) = 1 - P(B|A), \quad P[(B-C)|A] = P(B|A) - P(BC|A).
\]
\end{note}

\subsection{乘法公式}

\begin{theorem}[乘法公式] \label{thm:multiplication}
设 \(P(A) > 0\)，则
\[
P(AB) = P(B|A)\,P(A).
\]
一般地，设 \(A_1, A_2, \dots, A_n\) 为 \(n\) 个事件（\(n \ge 2\)），且 \(P(A_1 A_2 \cdots A_{n-1}) > 0\)，则
\[
P(A_1 A_2 \cdots A_n) = P(A_1)\,P(A_2|A_1)\,P(A_3|A_1 A_2) \cdots P(A_n|A_1 A_2 \cdots A_{n-1}).
\]
\end{theorem}

\subsection{全概率公式与贝叶斯公式}

前面已经介绍了完备事件组的概念（见定义~\ref{def:complete-event-group}），它是全概率公式和贝叶斯公式的前提条件。

\begin{theorem}[全概率公式] \label{thm:total-probability}
设 \(B_1, B_2, \dots, B_n\) 是一组完备事件组，\(A\) 为任一事件，且 \(P(B_i) > 0\ (i=1,2,\dots,n)\)，则
\[
P(A) = \sum_{i=1}^{n} P(B_i)\,P(A|B_i).
\]
\end{theorem}

全概率公式的实质是\textbf{全集分解法}：如果事件 \(A\) 的发生总是与多个"原因" \(B_i\) 相联系，则
\[
A = A\Omega = A\!\left(\bigcup_{i=1}^{n} B_i\right) = \bigcup_{i=1}^{n} B_i A,
\]
然后应用有限可加性和条件概率即可得到。

\begin{theorem}[贝叶斯公式] \label{thm:bayes}
设 \(B_1, B_2, \dots, B_n\) 是一组完备事件组，\(A\) 为任一事件，且 \(P(A) > 0\)，\(P(B_k) > 0\ (k=1,2,\dots,n)\)，则
\[
P(B_k|A) = \frac{P(B_k)\,P(A|B_k)}{P(A)} = \frac{P(B_k)\,P(A|B_k)}{\sum_{i=1}^{n} P(B_i)\,P(A|B_i)}, \quad k=1,2,\dots,n.
\]
\end{theorem}

\begin{note}
全概率公式与贝叶斯公式的记忆技巧：
\begin{itemize}
    \item \textbf{全概率公式}——由因求果。将 \(A\) 视为结果，\(B_i\) 视为导致 \(A\) 发生的原因，已知各原因的概率 \(P(B_i)\) 及各原因导致结果 \(A\) 的概率 \(P(A|B_i)\)，求结果 \(A\) 发生的概率 \(P(A)\)。
    \item \textbf{贝叶斯公式}——执果寻因。已知结果 \(A\) 已经发生，反推是由第 \(k\) 个原因 \(B_k\) 导致的概率 \(P(B_k|A)\)。
\end{itemize}
\end{note}

\begin{exercise}\label{exer:1-8}
某人外出可以乘坐飞机、火车、轮船、汽车四种交通工具，其概率依次为 \(0.05, 0.15, 0.30, 0.50\)，而乘坐这几种交通工具能如期到达的概率依次为 \(0.80, 0.70, 0.60, 0.90\)。求：
\par
(1) 该人如期到达的概率；\par
(2) 已知该人误期到达，求他是乘坐火车的概率。
\end{exercise}

\begin{solution}
设 \(A_i\ (i=1,2,3,4)\) 分别表示乘坐飞机、火车、轮船、汽车，\(B\) 表示如期到达。

(1) 利用全概率公式：
\[
P(B) = \sum_{i=1}^{4} P(A_i)\,P(B|A_i) = 0.05 \times 0.80 + 0.15 \times 0.70 + 0.30 \times 0.60 + 0.50 \times 0.90 = 0.775.
\]

(2) 利用贝叶斯公式：
\[
P(A_2|\overline{B}) = \frac{P(A_2)\,P(\overline{B}|A_2)}{P(\overline{B})} = \frac{0.15 \times (1-0.70)}{1-0.775} = \frac{0.045}{0.225} = 0.2.
\]
\end{solution}

\begin{exercise}\label{exer:1-9}
设 \(A, B\) 为两个随机事件，则（\quad）。\\
\begin{tabular}{ll}
A. \(P(AB)+P(\overline{A}\,\overline{B})\le 1\) & B. \(P(AB)+P(\overline{A}\,\overline{B})\ge 1\) \\[4pt]
C. \(P(A-B)\le P(A)-P(B)\) & D. \(P(A\cup B)+P(\overline{A}\cup\overline{B})\le 1\)
\end{tabular}
\end{exercise}

\begin{solution}
应选 A。

由于 \(AB \subset A \cup B\)，即 \(P(AB) \le P(A \cup B)\)，从而
\[
P(AB) \le P(A \cup B) = 1 - P(\overline{A \cup B}) = 1 - P(\overline{A}\,\overline{B}),
\]
即 \(P(AB) + P(\overline{A}\,\overline{B}) \le 1\)，故选 A。

同时有
\[
P(A \cup B) + P(\overline{A} \cup \overline{B}) - 1 = P(A \cup B) + P(\overline{A\,B}) - 1 \ge 0,
\]
故选项 B、D 不正确。又 \(AB \subset B\)，有 \(P(AB) \le P(B)\)，从而
\[
P(A-B) = P(A) - P(AB) \ge P(A) - P(B),
\]
知选项 C 不正确。
\end{solution}

\begin{exercise}\label{exer:1-10}
设随机事件 \(A, B\) 满足 \(P(A)=P(B)=\dfrac{1}{2}\) 和 \(P(A \cup B)=1\)，则（）。\par
A. \(A \cup B = \Omega\)\quad
B. \(AB = \varnothing\)\quad
C. \(P(\overline{A} \cup \overline{B})=1\)\quad
D. \(P(A-B)=0\)
\end{exercise}

\begin{solution}
应选 C。

\(P(A \cup B)=1\) 不能推出 \(A \cup B=\Omega\)，且 \(P(AB) = P(A \cup B) - P(A) - P(B) = 0\) 不能推出 \(AB=\varnothing\)。故选项 A、B 不正确。

又 \(P(A-B) = P(A) - P(AB) = \dfrac{1}{2}\)，选项 D 不正确。

从而 \(P(\overline{A} \cup \overline{B}) = P(\overline{AB}) = 1 - P(AB) = 1\)，故选 C。
\end{solution}

\begin{exercise}\label{exer:1-11}
设 \(A, B, C\) 是三个随机事件，\(A\) 与 \(C\) 互斥，\(P(AB)=\dfrac{1}{2}\)，\(P(C)=\dfrac{1}{3}\)，则 \(P(AB|\overline{C})=\underline{\hspace{2cm}}\)。
\end{exercise}

\begin{solution}
应填 \(\dfrac{3}{4}\)。

因为 \(A\) 与 \(C\) 互斥，所以 \(P(AC)=0\)。又 \(ABC \subset AC\)，所以 \(P(ABC)=0\)，故
\[
P(AB|\overline{C}) = \frac{P(AB\,\overline{C})}{P(\overline{C})} = \frac{P(AB) - P(ABC)}{1 - P(C)} = \frac{\frac{1}{2}}{1-\frac{1}{3}} = \frac{3}{4}.
\]
\end{solution}

\begin{note}
由于 \(AC = \varnothing\)，则 \(ABC = \varnothing \cdot B = \varnothing\)，得到 \(P(ABC)=0\)。
\end{note}

\begin{exercise}\label{exer:1-12}
已知 \(P(A)=\dfrac{1}{4}\)，\(P(B|A)=\dfrac{1}{3}\)，\(P(A|B)=\dfrac{1}{2}\)，则 \(P(A \cup B)=\underline{\hspace{2cm}}\)。
\end{exercise}

\begin{solution}
应填 \(\dfrac{1}{3}\)。

由乘法公式知
\[
P(AB) = P(A)\,P(B|A) = \frac{1}{4} \times \frac{1}{3} = \frac{1}{12}, \quad P(B) = \frac{P(AB)}{P(A|B)} = \frac{1/12}{1/2} = \frac{1}{6},
\]
所以
\[
P(A \cup B) = P(A) + P(B) - P(AB) = \frac{1}{4} + \frac{1}{6} - \frac{1}{12} = \frac{1}{3}.
\]
\end{solution}

\begin{exercise}\label{exer:1-13}
某人给四位亲友各写一封信，然后随机地装入 4 个写好地址的信封中，且每个信封装一封信。求：
\par
(1) 4 封信都装对的概率；\par
(2) 4 封信都装错的概率。
\end{exercise}

\begin{solution}
设事件 \(A_i\ (i=1,2,3,4)\) 为"第 \(i\) 封信装对了"，\(A\) 为"4 封信都装对了"，\(B\) 为"4 封信都装错了"。

(1) 利用乘法公式：
\[
\begin{aligned}
P(A) &= P(A_1 A_2 A_3 A_4) = P(A_1)\,P(A_2|A_1)\,P(A_3|A_1 A_2)\,P(A_4|A_1 A_2 A_3) \\
&= \frac{1}{4} \times \frac{1}{3} \times \frac{1}{2} \times \frac{1}{1} = \frac{1}{24}.
\end{aligned}
\]

(2) 利用逆向思维和对偶律：
\[
\begin{aligned}
P(B) &= P(\overline{A_1}\,\overline{A_2}\,\overline{A_3}\,\overline{A_4}) = 1 - P(A_1 \cup A_2 \cup A_3 \cup A_4) \\
&= 1 - \sum_{i=1}^{4} P(A_i) + \sum_{1 \le i < j \le 4} P(A_i A_j) - \sum_{1 \le i < j < k \le 4} P(A_i A_j A_k) + P(A_1 A_2 A_3 A_4),
\end{aligned}
\]
其中 \(P(A_i)=\dfrac{1}{4}\)，\(P(A_i A_j)=\dfrac{1}{4} \times \dfrac{1}{3}=\dfrac{1}{12}\)，\(P(A_i A_j A_k)=\dfrac{1}{4} \times \dfrac{1}{3} \times \dfrac{1}{2}=\dfrac{1}{24}\)，于是
\[
P(B) = 1 - 4 \times \frac{1}{4} + 6 \times \frac{1}{12} - 4 \times \frac{1}{24} + \frac{1}{24} = \frac{3}{8}.
\]
\end{solution}

\begin{note}
装信过程可看作无放回抽取信件，依次装入信封。第 (2) 问直接计算较为困难，采用逆向思维方法更为简捷。
\end{note}

\begin{exercise}\label{exer:1-14}
从 \(1, 2, 3, 4\) 中任取一个数，记为 \(X\)，再从 \(1, \dots, X\) 中任取一个数，记为 \(Y\)，则 \(P\{Y=2\}=\underline{\hspace{2cm}}\)。
\end{exercise}

\begin{solution}
应填 \(\dfrac{13}{48}\)。

由全概率公式：
\[
\begin{aligned}
P\{Y=2\} &= \sum_{x=1}^{4} P\{X=x\}\,P\{Y=2|X=x\} \\
&= \frac{1}{4} \times 0 + \frac{1}{4} \times \frac{1}{2} + \frac{1}{4} \times \frac{1}{3} + \frac{1}{4} \times \frac{1}{4} = \frac{13}{48}.
\end{aligned}
\]
\end{solution}


\begin{exercise}\label{exer:1-15}
一批产品有 10 个正品和 2 个次品，任意抽取两次，每次抽一个，抽出后不再放回，则第二次抽出的是次品的概率为 \(\underline{\hspace{2cm}}\)。
\end{exercise}

\begin{solution}
应填 \(\dfrac{1}{6}\)。

记 \(A_i=\{\text{第 }i\text{ 次取得次品}\}\ (i=1,2)\)，由全概率公式得
\[
P(A_2) = P(A_1)\,P(A_2|A_1) + P(\overline{A_1})\,P(A_2|\overline{A_1}) = \frac{2}{12} \times \frac{1}{11} + \frac{10}{12} \times \frac{2}{11} = \frac{1}{6}.
\]
\end{solution}

\begin{note}
此例为\textbf{抓阄模型}：无放回抽取时，每次取到次品的概率不变，均为 \(\dfrac{2}{12} = \dfrac{1}{6}\)。从第 1 次到第 12 次取到次品的概率均为 \(\dfrac{1}{6}\)。
\end{note}

\begin{exercise}\label{exer:1-16}
设有两批数量相同的零件，已知有一批产品全部合格，另一批产品有 25\% 不合格。从两批产品中任取 1 只，经检验是正品，放回原处，并在原所在批次再取 1 只，求这只产品是次品的概率。
\end{exercise}

\begin{solution}
设 \(H_i\ (i=1,2)\) 为"第一次从第 \(i\) 批产品中抽取"，\(A\) 为"所取产品为正品"。

第一次抽取：
\[
P(H_1)=P(H_2)=\frac{1}{2},\quad P(A|H_1)=1,\quad P(A|H_2)=\frac{3}{4},
\]
\[
P(A) = P(H_1)\,P(A|H_1) + P(H_2)\,P(A|H_2) = \frac{1}{2} \times 1 + \frac{1}{2} \times \frac{3}{4} = \frac{7}{8}.
\]

由贝叶斯公式更新概率：
\[
P(H_1|A) = \frac{P(H_1)\,P(A|H_1)}{P(A)} = \frac{4}{7}, \quad P(H_2|A) = \frac{3}{7}.
\]

第二次抽取：设所取为次品的概率为
\[
P(\overline{A}) = P(H_1|A) \cdot P(\overline{A}|H_1) + P(H_2|A) \cdot P(\overline{A}|H_2) = \frac{4}{7} \times 0 + \frac{3}{7} \times \frac{1}{4} = \frac{3}{28}.
\]
\end{solution}

\begin{note}
本题是全概率公式与贝叶斯公式的复合应用：第一次用全概率公式计算 \(P(A)\)，再用贝叶斯公式更新各原因的概率，最后第二次再用全概率公式计算结果。两次抽取的关键区别在于：第一次在完全不知情的情况下等可能地抽取；第二次在已知第一次取正品的条件下，各批次被抽到的概率已经发生了变化。
\end{note}

\begin{exercise}\label{exer:1-21}
设 \(P(\overline{A}) = 0.3\)，\(P(B) = 0.4\)，\(P(A\overline{B}) = 0.5\)，则 \(P(B|A\cup\overline{B}) = \underline{\hspace{2cm}}\)。
\end{exercise}

\begin{solution}
应填 \(0.25\)。

首先，\(P(A) = 1 - P(\overline{A}) = 0.7\)。根据减法公式：
\[P(AB) = P(A) - P(A\overline{B}) = 0.7 - 0.5 = 0.2.\]
根据加法公式：
\[P(A\cup\overline{B}) = P(A) + P(\overline{B}) - P(A\overline{B}) = 0.7 + 0.6 - 0.5 = 0.8.\]
利用分配律化简分子：
\[B \cap (A\cup\overline{B}) = (B\cap A)\cup(B\cap\overline{B}) = AB,\]
故
\[P(B|A\cup\overline{B}) = \frac{P(AB)}{P(A\cup\overline{B})} = \frac{0.2}{0.8} = 0.25.\]
\end{solution}

\begin{exercise}\label{exer:1-22}
设 \(A, B\) 为任意两事件且 \(A\subset B\)，\(P(B) > 0\)，则下列结论一定正确的是（\quad）。\par
A. \(P(A) \ge P(A|B)\)\quad
B. \(P(A) > P(A|B)\)\quad
C. \(P(A) < P(A|B)\)\quad
D. \(P(A) \le P(A|B)\)
\end{exercise}

\begin{solution}
应选 D。

由 \(A\subset B\)，有 \(AB = A\)，故
\[P(A|B) = \frac{P(AB)}{P(B)} = \frac{P(A)}{P(B)} \ge P(A).\]

从样本空间更新的角度理解：\(P(A) = P(A|\Omega)\)，而 \(A\subset B\subset\Omega\)，样本空间缩小为 \(B\) 后，\(A\) 在新样本空间中的"占比"不变，但不确定性减小，概率增大。
\end{solution}

\begin{exercise}\label{exer:1-23}
假如在数字通信中传送信号 0 与 1 的概率分别为 0.7 和 0.3。由于随机干扰，当传送信号 0 时接收到信号 0 的概率是 0.8；当传送信号 1 时接收到信号 1 的概率是 0.9。求：\par
(1) 接收到信号 0 的概率；\par
(2) 当接收到信号 0 时，传送的信号是 0 的概率。
\end{exercise}

\begin{solution}
设事件 \(A\)：传送信号 0；事件 \(B\)：接收到信号 0。已知 \(P(A) = 0.7\)，\(P(\overline{A}) = 0.3\)，\(P(B|A) = 0.8\)，\(P(\overline{B}|\overline{A}) = 0.9\)，故 \(P(B|\overline{A}) = 0.1\)。

(1) 由全概率公式：
\[P(B) = P(B|A)\,P(A) + P(B|\overline{A})\,P(\overline{A}) = 0.8 \times 0.7 + 0.1 \times 0.3 = 0.59.\]

(2) 由贝叶斯公式：
\[P(A|B) = \frac{P(B|A)\,P(A)}{P(B)} = \frac{0.8 \times 0.7}{0.59} = \frac{56}{59}.\]
\end{solution}

\begin{note}
本题是全概率公式与贝叶斯公式的典型应用：传送信号为"因"，接收信号为"果"。第 (1) 问"由因求果"用全概率公式，第 (2) 问"执果寻因"用贝叶斯公式。
\end{note}

\begin{exercise}\label{exer:1-39}
甲、乙两人独立地对同一目标射击一次，其命中率分别为 0.6 和 0.7。已知目标被击中，则它是甲射中的概率为 \underline{\hspace{2cm}}。
\end{exercise}

\begin{solution}
应填 \(\dfrac{15}{22}\)（约 0.682）。

记 \(A\) 为"目标被击中"，\(B_1\) 为"甲射中"，\(B_2\) 为"乙射中"。\(A = B_1\cup B_2\)，由独立性：
\[P(A) = P(B_1)+P(B_2)-P(B_1B_2) = 0.6+0.7-0.6\times 0.7 = 0.88.\]
由于 \(B_1\subset A\)，
\[P(B_1|A) = \frac{P(B_1)}{P(A)} = \frac{0.6}{0.88} = \frac{15}{22}.\]
\end{solution}

\begin{exercise}\label{exer:1-40}
口袋中有 1 个球，不知其颜色是黑还是白。现再往口袋中放入 1 个白球，然后从中任意取出 1 个，发现取出的是白球。求口袋中原来那个球是白球的概率。
\end{exercise}

\begin{solution}
应填 \(\dfrac{2}{3}\)。

记 \(A\) 为"取出的是白球"，\(B\) 为"原来那个球是白球"。由于不知道原球颜色，取等可能先验概率 \(P(B)=P(\overline{B})=\dfrac{1}{2}\)。

\(P(A|B)=1\)（原来若为白球，袋中有 2 白，必取白），\(P(A|\overline{B})=\dfrac{1}{2}\)（原来若为黑球，袋中 1 白 1 黑）。

由贝叶斯公式：
\[P(B|A) = \frac{P(B)\,P(A|B)}{P(B)\,P(A|B)+P(\overline{B})\,P(A|\overline{B})} = \frac{\frac{1}{2}\times 1}{\frac{1}{2}\times 1+\frac{1}{2}\times\frac{1}{2}} = \frac{2}{3}.\]
\end{solution}

\begin{exercise}\label{exer:1-41}
已知装有同种零件的产品两箱，第一箱内装 50 件，其中一等品 10 件；第二箱内装 30 件，其中一等品 18 件。现从两箱中任意挑选一箱，从中先后取出两件产品（取后不放回）。求：先取出的产品是一等品的概率 \(p\)；已知先取出的产品是一等品，则第二次取出的产品仍然是一等品的概率 \(q\)。
\end{exercise}

\begin{solution}
设 \(A_i\) 为"第 \(i\) 次取出一等品"（\(i=1,2\)），\(C_j\) 为"挑选第 \(j\) 箱"（\(j=1,2\)）。\(C_1\cup C_2=\Omega\)，\(C_1C_2=\varnothing\)，\(P(C_1)=P(C_2)=\dfrac{1}{2}\)。

由全概率公式：
\begin{align*}
p = P(A_1) &= P(C_1)\,P(A_1|C_1)+P(C_2)\,P(A_1|C_2) = \frac{1}{2}\times\frac{10}{50}+\frac{1}{2}\times\frac{18}{30} = \frac{2}{5}, \\
P(A_1A_2) &= P(C_1)\,P(A_1A_2|C_1)+P(C_2)\,P(A_1A_2|C_2) \\
&= \frac{1}{2}\times\frac{10\times 9}{50\times 49}+\frac{1}{2}\times\frac{18\times 17}{30\times 29} = \frac{276}{1421}.
\end{align*}
故 \(q = P(A_2|A_1) = \dfrac{P(A_1A_2)}{P(A_1)} = \dfrac{276}{1421}\times\dfrac{5}{2} = \dfrac{690}{1421}\)。
\end{solution}

\begin{exercise}\label{exer:1-42}
每箱产品有 10 件，其次品数从 0 到 2 是等可能的。开箱检验时从中任取 1 件，检验出次品则拒收。由于检验有误，将正品误认为次品的概率为 2\%，次品被漏查的概率为 5\%。求该箱产品通过验收的概率。
\end{exercise}

\begin{solution}
应填 \(0.887\)。

设 \(A_i\)（\(i=0,1,2\)）为"箱中有 \(i\) 件次品"，\(B\) 为"通过验收"。\(P(A_i)=\dfrac{1}{3}\)。

取到正品时通过验收的概率为 0.98，取到次品时通过验收的概率为 0.05。由全概率公式（两层）：
\begin{align*}
P(B|A_0) &= 1\times 0.98 = 0.98, \\
P(B|A_1) &= \frac{9}{10}\times 0.98 + \frac{1}{10}\times 0.05 = 0.887, \\
P(B|A_2) &= \frac{8}{10}\times 0.98 + \frac{2}{10}\times 0.05 = 0.794, \\
P(B) &= \frac{1}{3}(0.98+0.887+0.794) = \frac{2.661}{3} = 0.887.
\end{align*}
\end{solution}

\begin{exercise}\label{exer:1-52}
设水杯成箱出售，每箱 12 个，每箱中含 0 个、1 个、2 个残次品的概率分别为 0.7、0.19、0.11。某人欲购一箱，售货员随机取一箱，从中任取 3 只查看，若无残次品则购买。\par
(1) 此人购买一箱的概率；\par
(2) 在此人购买一箱的条件下，该箱无残次品的概率。
\end{exercise}

\begin{solution}
设 \(A_i\)（\(i=0,1,2\)）为"箱中有 \(i\) 件残次品"，\(B\) 为"购买此箱"（即取出的 3 只均无残次品）。

(1) 由全概率公式：
\begin{align*}
P(B) &= \sum_{i=0}^{2}P(A_i)\,P(B|A_i), \\
P(B|A_0) &= 1, \quad P(B|A_1) = \frac{\mathrm{C}_{11}^3}{\mathrm{C}_{12}^3} = \frac{3}{4}, \quad P(B|A_2) = \frac{\mathrm{C}_{10}^3}{\mathrm{C}_{12}^3} = \frac{6}{11}, \\
P(B) &= 0.7\times 1 + 0.19\times\frac{3}{4} + 0.11\times\frac{6}{11} = 0.9025.
\end{align*}

(2) 由贝叶斯公式：
\[
P(A_0|B) = \frac{P(A_0)\,P(B|A_0)}{P(B)} = \frac{0.7\times 1}{0.9025} = \frac{280}{361} \approx 0.7756.
\]
\end{solution}

\begin{exercise}\label{exer:1-53}
统计资料显示，在出口罐头导致的索赔事件中，有 50\% 是质量问题，30\% 是数量问题，20\% 是包装问题。在这三类问题各引起的索赔事件中，不经过法律诉讼而通过协商解决的概率分别为 40\%、60\% 和 75\%。\par
(1) 索赔事件通过协商解决的概率；\par
(2) 已知一索赔事件通过协商解决，求该事件不是包装问题的概率。
\end{exercise}

\begin{solution}
设 \(A_1,A_2,A_3\) 分别为质量、数量、包装问题引起的索赔，\(B\) 为"通过协商解决"。

(1) 由全概率公式：
\[
P(B) = \sum_{i=1}^{3}P(A_i)\,P(B|A_i) = 0.5\times 0.40 + 0.3\times 0.60 + 0.2\times 0.75 = 0.53.
\]

(2) 由贝叶斯公式：
\[
P(A_3|B) = \frac{P(A_3)\,P(B|A_3)}{P(B)} = \frac{0.2\times 0.75}{0.53} = \frac{15}{53},
\]
故所求概率为 \(1-P(A_3|B) = \dfrac{38}{53}\)。
\end{solution}

\begin{exercise}\label{exer:1-54}
在某城市中，下雨和晴天的天数各占一半，而天气无论在雨天还是晴天都有 \(\dfrac{2}{3}\) 的概率预报准确。如果预报下雨，王明同学就一定带雨伞；如果预报无雨，则不带伞。设 \(A=\{\text{天下雨}\}\)，\(B=\{\text{预报有雨}\}\)，\(C=\{\text{王明带雨伞}\}\)。\par
(1) 事件 \(\overline{A}\cap\overline{B}\cap C\)、\(\overline{A}\cap B\cap C\) 的含义是什么？哪个为不可能事件？\par
(2) 求他带雨伞而没有下雨的概率。
\end{exercise}

\begin{solution}
由题意：\(P(A)=\dfrac{1}{2}\)，\(P(B|A)=\dfrac{2}{3}\)（预报准确），\(P(\overline{B}|\overline{A})=\dfrac{2}{3}\)（预报准确），故 \(P(B|\overline{A})=\dfrac{1}{3}\)。又 \(P(C|B)=1\)（预报有雨必带伞），\(P(C|\overline{B})=0\)（预报无雨不带伞）。

(1) \(\overline{A}\cap\overline{B}\cap C\)：天没下雨、预报无雨、却带了雨伞。由于 \(P(C|\overline{B})=0\)，这是一个\textbf{不可能事件}。\par
\(\overline{A}\cap B\cap C\)：天没下雨、预报有雨、带了雨伞。这是可能事件。

(2) "带雨伞而没有下雨"即 \(\overline{A}\cap C\)。由于 \(C\) 完全由 \(B\) 决定，\(\overline{A}\cap\overline{B}\cap C\) 为不可能事件，故
\[
P(\overline{A}\cap C) = P(\overline{A}\cap B\cap C) = P(\overline{A})\,P(B|\overline{A})\,P(C|\overline{A}B) = \frac{1}{2}\times\frac{1}{3}\times 1 = \frac{1}{6}.
\]
\end{solution}


\begin{exercise}[提高]\label{exer:1-55}
设某同学的手机在一天内收到短信数服从泊松分布 \(P(\lambda)\)，每个短信是否为垃圾短信与其到达时间独立，也与其他短信是否为垃圾短信相互独立。假设每个短信是垃圾短信的概率为 \(p\)。\par
(1) 已知一天内收到了 \(n\) 条短信，求其中恰有 \(k\) 条垃圾短信的概率（\(0\le k\le n\)）。\par
(2) 求一天内收到 \(k\) 条垃圾短信的概率（\(k=0,1,2,\dots\)）。
\end{exercise}

\begin{solution}
(1) 记 \(X\) 为一天内收到短信的总数，\(Y\) 为其中垃圾短信的条数。已知 \(X=n\) 时，每条短信独立地以概率 \(p\) 为垃圾短信，故
\[
P(Y=k\mid X=n) = \mathrm{C}_n^k\,p^k\,(1-p)^{n-k}.
\]

(2) 由全概率公式：
\begin{align*}
P(Y=k) &= \sum_{n=k}^{\infty} P(Y=k\mid X=n)\,P(X=n) \\
&= \sum_{n=k}^{\infty} \mathrm{C}_n^k\,p^k\,(1-p)^{n-k}\,\frac{\lambda^n}{n!}\,e^{-\lambda} \\
&= \frac{p^k\,e^{-\lambda}\,\lambda^k}{k!}\sum_{n=k}^{\infty}\frac{[\lambda(1-p)]^{n-k}}{(n-k)!} \\
&= \frac{(\lambda p)^k}{k!}\,e^{-\lambda}\,e^{\lambda(1-p)} = \frac{(\lambda p)^k}{k!}\,e^{-\lambda p}.
\end{align*}
即 \(Y\sim P(\lambda p)\)。
\end{solution}

\begin{note}
泊松分布的分解定理（稀疏性）：以概率 \(p\) 独立地对泊松过程 \(P(\lambda)\) 进行"稀疏"，得到的新过程仍为泊松分布，参数变为 \(\lambda p\)。本题是全概率公式与泊松分布结合的典型应用。
\end{note}


\section{事件的独立性}

\subsection{独立性的定义}

\begin{definition}[事件的独立性] \label{def:independence}
设 \(A, B\) 为两个事件，如果
\[
P(AB) = P(A)\,P(B),
\]
则称事件 \(A\) 与 \(B\) \textbf{相互独立}，简称 \(A\) 与 \(B\) 独立。
\end{definition}

\begin{definition}[$n$ 个事件的相互独立] \label{def:n-independence}
设 \(A_1, A_2, \dots, A_n\) 为 \(n\ (n \ge 2)\) 个事件，如果对其中任意有限个事件 \(A_{i_1}, A_{i_2}, \dots, A_{i_k}\ (2 \le k \le n)\)，都有
\[
P(A_{i_1} A_{i_2} \cdots A_{i_k}) = P(A_{i_1})\,P(A_{i_2}) \cdots P(A_{i_k}),
\]
则称这 \(n\) 个事件\textbf{相互独立}。
\end{definition}

特别地，对于三个事件 \(A, B, C\)，相互独立要求同时满足以下四个等式：
\[
\begin{cases}
P(AB) = P(A)\,P(B), \\
P(AC) = P(A)\,P(C), \\
P(BC) = P(B)\,P(C), \\
P(ABC) = P(A)\,P(B)\,P(C).
\end{cases}
\]
若仅满足前三个等式（即任意两个事件独立），则称 \(A, B, C\) \textbf{两两独立}。

\begin{note}
从条件概率的角度理解独立性：若 \(A, B\) 相互独立且 \(P(B) > 0\)，则 \(P(A|B) = P(A)\)。这意味着样本空间缩小为 \(B\) 后，事件 \(A\) 发生的概率并没有改变——\(B\) 的发生对 \(A\) 没有任何影响，这正是"独立"一词的含义。
\end{note}

\subsection{独立性的判定}

\begin{property}\label{prop:independence-equiv}
\(A\) 与 \(B\) 独立（\(P(AB)=P(A)\,P(B)\)）的等价条件（设 \(0 < P(A) < 1\)）：
\begin{enumerate}
    \item \(A\) 与 \(B\) 独立 \(\Leftrightarrow\) \(A\) 与 \(\overline{B}\) 独立 \(\Leftrightarrow\) \(\overline{A}\) 与 \(B\) 独立 \(\Leftrightarrow\) \(\overline{A}\) 与 \(\overline{B}\) 独立；
    \item \(A\) 与 \(B\) 独立 \(\Leftrightarrow\) \(P(B|A) = P(B|\overline{A}) = P(B)\)；
    \item \(A\) 与 \(B\) 独立 \(\Leftrightarrow\) \(P(B|A) + P(\overline{B}|\overline{A}) = 1\)。
\end{enumerate}
\end{property}

\begin{remark}
(1) 关于等价条件 (1) 的证明：由 \(A, B\) 独立，有 \(P(AB)=P(A)\,P(B)\)，于是
\[
P(A\,\overline{B}) = P(A) - P(AB) = P(A) - P(A)\,P(B) = P(A)\,[1-P(B)] = P(A)\,P(\overline{B}),
\]
故 \(A, \overline{B}\) 独立，其余同理。\newline
(2) 推广：将相互独立的事件组中的任何几个事件换成各自的对立事件，所得的新事件组仍相互独立。
\end{remark}

\begin{property}\label{prop:independence-compute}
独立性的判定与运算：
\begin{enumerate}
    \item 若 \(P(A)=0\) 或 \(P(A)=1\)，则 \(A\) 与任意事件 \(B\) 相互独立。
    \item 若 \(P(A)>0, P(B)>0\)，且 \(A, B\) 互斥，则 \(A, B\) 一定\textbf{不}独立；反之，若 \(A, B\) 独立，则 \(A, B\) 一定\textbf{不}互斥（除非 \(P(A)=0\) 或 \(P(B)=0\)）。
    \item 若一组事件相互独立，则其中部分事件也相互独立；对不含相同事件的独立组分别作运算，所得新事件仍独立。例如 \(A, B, C, D\) 相互独立，则 \(AB\) 与 \(CD\) 独立，\(A\) 与 \(BC-D\) 独立。
\end{enumerate}
\end{property}

\begin{note}
常见误区：
\begin{enumerate}
    \item "独立性"是概率意义上的数学定义，与逻辑因果上的"独立"无关，只看 \(P(AB)=P(A)\,P(B)\) 是否成立。
    \item 两两独立 \(\nRightarrow\) 相互独立。三个事件两两独立只保证 \(P(AB)=P(A)\,P(B)\)、\(P(AC)=P(A)\,P(C)\)、\(P(BC)=P(B)\,P(C)\)，但不保证 \(P(ABC)=P(A)\,P(B)\,P(C)\)。
\end{enumerate}
\end{note}

\begin{exercise}\label{exer:1-17}
将一枚硬币独立地掷两次，设 \(A_1=\{\text{掷第一次出现正面}\}\)，\(A_2=\{\text{掷第二次出现正面}\}\)，\(A_3=\{\text{正反面各出现一次}\}\)，\(A_4=\{\text{正面出现两次}\}\)，则事件（\quad）。\par
\begin{tabular}{ll}
A. \(A_1,A_2,A_3\) 相互独立 & B. \(A_2,A_3,A_4\) 相互独立 \\[4pt]
C. \(A_1,A_2,A_3\) 两两独立 & D. \(A_2,A_3,A_4\) 两两独立
\end{tabular}
\end{exercise}

\begin{solution}
应选 C。

根据古典概型：
\[
P(A_1)=P(A_2)=P(A_3)=\frac{1}{2}, \quad P(A_1A_2)=P(A_1A_3)=P(A_2A_3)=\frac{1}{4},
\]
满足两两独立条件。但
\[
P(A_1A_2A_3)=0 \neq P(A_1)\,P(A_2)\,P(A_3)=\frac{1}{8},
\]
故 \(A_1, A_2, A_3\) 两两独立但不相互独立。

又 \(A_3 \cap A_4 = \varnothing\)，\(P(A_3A_4)=0 \neq P(A_3)\,P(A_4)=\frac{1}{8}\)，故 \(A_2, A_3, A_4\) 不两两独立。
\end{solution}

\begin{exercise}\label{exer:1-18}
设 \(A, B, C\) 为三个两两独立的随机事件，\(P(A)=P(B)=P(C)\)，\(P(A \cup B \cup C)=\dfrac{23}{25}\)，\(P(\overline{A} \cup \overline{B} \cup \overline{C})=\dfrac{14}{25}\)，求 \(P(A)\)。
\end{exercise}

\begin{solution}
由两两独立性，
\[
\begin{aligned}
P(\overline{A} \cup \overline{B} \cup \overline{C}) &= P(\overline{A})+P(\overline{B})+P(\overline{C}) - P(\overline{A}\,\overline{B}) - P(\overline{B}\,\overline{C}) - P(\overline{A}\,\overline{C}) + P(\overline{A}\,\overline{B}\,\overline{C}) \\
&= 3P(\overline{A}) - 3[P(\overline{A})]^2 + P(A \cup B \cup C).
\end{aligned}
\]
代入数据：
\[
3P(\overline{A}) - 3[P(\overline{A})]^2 + \frac{2}{25} = \frac{14}{25} \implies 3P(\overline{A}) - 3[P(\overline{A})]^2 = \frac{12}{25}.
\]
令 \(x = P(\overline{A})\)，则 \(3x - 3x^2 = \dfrac{12}{25}\)，即 \(25x^2 - 25x + 4 = 0\)，解得 \(x = \dfrac{1}{5}\) 或 \(\dfrac{4}{5}\)。

又 \(P(\overline{A}) \le P(\overline{A} \cup \overline{B} \cup \overline{C}) = \dfrac{14}{25}\)，故 \(P(\overline{A}) = \dfrac{1}{5}\)，即 \(P(A) = \dfrac{4}{5}\)。
\end{solution}

\begin{exercise}\label{exer:1-43}
设 \(A,B,C\) 为三个随机事件，且 \(A\) 与 \(C\) 相互独立，\(B\) 与 \(C\) 相互独立，则 \(A\cup B\) 与 \(C\) 相互独立的充分必要条件是（\quad）。\par
A. \(A\) 与 \(B\) 相互独立\quad
B. \(A\) 与 \(B\) 互不相容\quad
C. \(AB\) 与 \(C\) 相互独立\quad
D. \(AB\) 与 \(C\) 互不相容
\end{exercise}

\begin{solution}
应选 C。

\(A\cup B\) 与 \(C\) 相互独立，即 \(P[(A\cup B)C] = P(A\cup B)\,P(C)\)。
\[
\begin{aligned}
P[(A\cup B)C] &= P(AC\cup BC) = P(AC)+P(BC)-P(ABC) \\
&= P(A)\,P(C)+P(B)\,P(C)-P(ABC), \\
P(A\cup B)\,P(C) &= [P(A)+P(B)-P(AB)]\,P(C) \\
&= P(A)\,P(C)+P(B)\,P(C)-P(AB)\,P(C),
\end{aligned}
\]
因此充要条件为 \(P(ABC)=P(AB)\,P(C)\)，即 \(AB\) 与 \(C\) 相互独立。
\end{solution}

\begin{exercise}\label{exer:1-56}
已知 \(A,B\) 相互独立，\(P(A)=0.5\)，\(P(B)=0.4\)，\(C=\overline{A}\cup B\)。求 \(P(\overline{C})\)，\(P(C\mid(A\cup B))\)。
\end{exercise}

\begin{solution}
\(\overline{C}=\overline{\overline{A}\cup B}=A\,\overline{B}\)，由独立性：
\[
P(\overline{C}) = P(A\,\overline{B}) = P(A)(1-P(B)) = 0.5\times 0.6 = 0.3.
\]

又 \(P(A\cup B)=P(A)+P(B)-P(A)P(B)=0.5+0.4-0.2=0.7\)。注意到
\[
C\cap(A\cup B) = (\overline{A}\cup B)\cap(A\cup B) = B,
\]
故
\[
P(C\mid(A\cup B)) = \frac{P(B)}{P(A\cup B)} = \frac{0.4}{0.7} = \frac{4}{7}.
\]
\end{solution}

\begin{exercise}\label{exer:1-57}
下列命题中，正确的是（\quad）。\par
A. 若 \(A,B\) 互不相容，则 \(\overline{A},\overline{B}\) 也互不相容\quad
B. 若 \(A,B\) 相容，则 \(\overline{A},\overline{B}\) 也相容\quad
C. 若 \(A,B\) 独立，则 \(\overline{A},\overline{B}\) 也独立\quad
D. 若 \(A,B\) 独立，则 \(\overline{A},\overline{B}\) 不相容
\end{exercise}

\begin{solution}
应选 C。

由独立性 \(P(AB)=P(A)P(B)\)：
\[
P(\overline{A}\,\overline{B}) = P(\overline{A\cup B}) = 1-P(A)-P(B)+P(AB) = (1-P(A))(1-P(B)) = P(\overline{A})P(\overline{B}),
\]
故 \(\overline{A},\overline{B}\) 也独立。

A 错误：例如 \(A=\varnothing\) 时 \(\overline{A}=\Omega\)，\(\overline{A}\,\overline{B}=\overline{B}\) 一般不为空。B 错误：类似地，\(\overline{A},\overline{B}\) 可相容也可不相容。D 错误：独立与相容没有必然关系。
\end{solution}


\begin{exercise}[提高]\label{exer:1-58}
设 \(A,B,C\) 相互独立，且 \(P(A)\neq 0\)，\(0<P(C)<1\)。则下面四个事件中不独立的是（\quad）。\par
A. \(\overline{A}\cup B\) 与 \(C\)\quad
B. \(\overline{AC}\) 与 \(\overline{C}\)\quad
C. \(\overline{A-B}\) 与 \(\overline{C}\)\quad
D. \(\overline{AB}\) 与 \(\overline{C}\)
\end{exercise}

\begin{solution}
应选 B。

A、C、D 中左边的事件均由 \(A,B\) 构成，而 \(A,B\) 各自与 \(C\) 独立，因此由 \(A,B\) 经过并、交、差运算构成的新事件也与 \(C\) 独立。

B 中，\(\overline{AC}\cap\overline{C}=\overline{C}\)（因为 \(\overline{AC}\supset\overline{C}\)）。若两者独立，需 \(P(\overline{AC})\,P(\overline{C})=P(\overline{C})\)，即 \(P(\overline{AC})=1\)。但
\[
P(\overline{AC}) = P(\overline{A})+P(\overline{C})-P(\overline{A})P(\overline{C}) < 1 \quad (\text{因为 }P(\overline{A})<1, P(\overline{C})<1),
\]
矛盾。故 \(\overline{AC}\) 与 \(\overline{C}\) 不独立。

等价地，若独立则需 \(P(\overline{C})=P(\overline{AC}\,\overline{C})=P(\overline{C})\) 恒成立，但由独立性 \(P(AC)=P(A)P(C)>0\)，可得 \(\overline{AC}\neq\Omega\)，从而 \(P(\overline{AC})<1\)。
\end{solution}

\section{独立重复试验与伯努利概型}

\begin{definition}[\(n\) 重伯努利概型] \label{def:bernoulli}
如果随机试验满足以下条件：
\begin{enumerate}
    \item 每次试验只有两个可能结果，记为 \(A\)（成功）和 \(\overline{A}\)（失败）；
    \item 每次试验中 \(P(A)=p\)（\(0 < p < 1\)）保持不变；
    \item 各次试验相互独立，
\end{enumerate}
则将这种试验独立重复进行 \(n\) 次的概率模型称为 \(\boldsymbol{n}\) \textbf{重伯努利概型}。
\end{definition}

在 \(n\) 重伯努利概型中，事件 \(A\) 恰好发生 \(k\) 次的概率为
\[
P\{A \text{ 恰好发生 } k \text{ 次}\} = \mathrm{C}_n^k\, p^k\, (1-p)^{n-k}, \quad k=0,1,\dots,n.
\]

如果用 \(X\) 表示 \(n\) 重伯努利概型中事件 \(A\) 发生的次数，则 \(X\) 服从参数为 \(n, p\) 的\textbf{二项分布}，记为 \(X \sim B(n,p)\)。

\begin{note}
要善于判定 \(n\) 重伯努利概型：题目中出现“将……独立重复进行 \(n\) 次”“对……重复观察 \(n\) 次”等字样，或可以转化为 \(n\) 次独立重复试验的问题，都应考虑使用二项分布计算概率。例如射击问题、抛硬币问题、产品抽样检验问题等。
\end{note}

\begin{exercise}\label{exer:1-19}
对同一目标接连进行 3 次独立重复射击，假设至少命中目标一次的概率为 \(\dfrac{7}{8}\)，则每次射击命中目标的概率 \(p=\underline{\hspace{2cm}}\)。
\end{exercise}

\begin{solution}
应填 \(\dfrac{1}{2}\)。

记 \(A_i=\{\text{第 }i\text{ 次命中}\}\ (i=1,2,3)\)，\(A_1, A_2, A_3\) 相互独立且 \(P(A_i)=p\)。由
\[
P(A_1 \cup A_2 \cup A_3) = 1 - P(\overline{A_1}\,\overline{A_2}\,\overline{A_3}) = 1 - (1-p)^3 = \frac{7}{8},
\]
得 \((1-p)^3 = \dfrac{1}{8}\)，故 \(1-p = \dfrac{1}{2}\)，\(p = \dfrac{1}{2}\)。
\end{solution}

\begin{exercise}\label{exer:1-20}
某人向同一目标独立重复射击，每次命中目标的概率为 \(p\ (0<p<1)\)，则此人第 4 次射击恰好第 2 次命中目标的概率为（\quad）。\par
A. \(3p(1-p)^2\)\quad
B. \(6p(1-p)^2\)\quad
C. \(3p^2(1-p)^2\)\quad
D. \(6p^2(1-p)^2\)
\end{exercise}

\begin{solution}
应选 C。

"第 4 次恰好第 2 次命中"意味着：前 3 次中有 1 次命中，第 4 次命中。设前 3 次命中次数为 \(X \sim B(3,p)\)，则
\[
P\{X=1\} = \mathrm{C}_3^1\, p\,(1-p)^2 = 3p(1-p)^2.
\]
第 4 次命中的概率为 \(p\)，故所求概率为
\[
3p(1-p)^2 \cdot p = 3p^2(1-p)^2.
\]
\end{solution}

\begin{note}
一般地，在独立重复试验中，第 \(n\) 次试验恰好第 \(k\) 次成功的概率为
\[
\mathrm{C}_{n-1}^{k-1}\, p^k\, (1-p)^{n-k},
\]
其中前 \(n-1\) 次中有 \(k-1\) 次成功，第 \(n\) 次成功。这实际上是负二项分布（帕斯卡分布）的一个特例。
\end{note}

\begin{exercise}\label{exer:1-25}
在四次独立试验中，事件 \(A\) 至少发生一次的概率为 \(\dfrac{80}{81}\)，则事件 \(A\) 在每次试验中发生的概率为（\quad）。\par
A. \(\dfrac{1}{3}\)\quad
B. \(\dfrac{2}{3}\)\quad
C. \(\dfrac{2}{3}\sqrt[4]{5}\)\quad
D. \(1 - \dfrac{2}{3}\sqrt[4]{5}\)
\end{exercise}

\begin{solution}
应选 B。设 \(X\) 为事件 \(A\) 在 4 次试验中发生的次数。
\[P\{X \ge 1\} = 1 - P\{X = 0\} = 1 - (1-p)^4 = \frac{80}{81},\]
故 \((1-p)^4 = \dfrac{1}{81}\)，即 \(1-p = \dfrac{1}{3}\)，\(p = \dfrac{2}{3}\)。
\end{solution}

\begin{exercise}\label{exer:1-44}
三人独立地破译一个密码，他们能单独译出的概率分别为 \(\dfrac{1}{5},\dfrac{1}{3},\dfrac{1}{4}\)，则此密码被译出的概率为 \underline{\hspace{2cm}}。
\end{exercise}

\begin{solution}
应填 \(\dfrac{3}{5}\)。

记事件 \(A_i\) 为"第 \(i\) 个人译出密码"（\(i=1,2,3\)），\(B\) 为"密码被译出"。由独立性：
\[P(B) = P(A_1\cup A_2\cup A_3) = 1-P(\overline{A_1}\,\overline{A_2}\,\overline{A_3}) = 1-\frac{4}{5}\times\frac{2}{3}\times\frac{3}{4} = \frac{3}{5}.\]
\end{solution}

\begin{note}
互不相容可简化并事件的概率计算，相互独立可简化交事件的概率计算。这里利用对偶律将事件的并转化为事件的交，从而利用独立性进行计算，这一方法经常用到。
\end{note}

\begin{exercise}\label{exer:1-59}
设甲、乙两人轮流射击同一目标，直到射中为止，每次命中率分别为 0.3 和 0.4，甲先射击。\par
(1) 目标是由甲射中的概率 \(p_1\)；\par
(2) 击中目标者射击了 \(k\) 次的概率 \(p_2\)。
\end{exercise}

\begin{solution}
(1) 甲在第一轮命中的概率为 0.3；若第一轮都未命中（概率 \(0.7\times 0.6=0.42\)），则回到初始状态。故
\[
p_1 = 0.3 + 0.42\times p_1 \implies p_1 = \frac{0.3}{1-0.42} = \frac{15}{29} \approx 0.5172.
\]

(2) "击中者射击了 \(k\) 次"意味着该射手在前 \(k-1\) 轮均未命中，第 \(k\) 次命中。若甲击中：前 \(k-1\) 轮甲乙均未命中（概率 \(0.42^{k-1}\)），甲第 \(k\) 次命中（概率 0.3）。若乙击中：前 \(k-1\) 轮均未命中（概率 \(0.42^{k-1}\)），甲第 \(k\) 次未命中（概率 0.7），乙第 \(k\) 次命中（概率 0.4）。故
\[
p_2 = 0.42^{k-1}\times 0.3 + 0.42^{k-1}\times 0.7\times 0.4 = 0.42^{k-1}\times 0.58 \quad (k=1,2,\dots).
\]
\end{solution}

\chapter{一维随机变量及其分布}

\[
\begin{cases}
\text{随机变量分布函数 } F(x) = P(X \le x) \text{ 的性质}
\begin{cases}
(1)\ \text{单调不减} \\
(2)\ 0 \le F(x) \le 1,\ F(-\infty)=0,\ F(+\infty)=1 \\
(3)\ \text{右连续: } F(x+0) = F(x)
\end{cases} \\[1em]
\text{随机变量分类及概率分布}
\begin{cases}
\text{离散型: 分布律} \\
\text{连续型: 概率密度} \\
\text{混合型: 分布函数}
\end{cases}
\ \text{满足概率分布的性质} \\[1em]
\text{常见随机变量及其分布}
\begin{cases}
\text{离散型}
\begin{cases}
(0-1)\text{分布: } P(X=k) = p^k(1-p)^{1-k}\ (k=0,1) \\
\text{二项分布: } P(X=k) = C_n^k p^k(1-p)^{n-k}\ (k=0,1,2,\dots,n) \\
\text{泊松分布: } P(X=k) = \frac{\lambda^k}{k!}e^{-\lambda}\ (k=0,1,2\cdots)
\end{cases} \\[2em]
\text{连续型}
\begin{cases}
\text{均匀分布}
\begin{cases}
X \sim U(a,b) \\
\text{性质: 区间概率等于该区间占整个区间的比例}
\end{cases} \\
\text{指数分布}
\begin{cases}
X \sim E(\lambda) \\
\text{性质: 无记忆性}
\end{cases} \\
\text{正态分布}
\begin{cases}
X \sim N(\mu,\sigma^2) \\
\text{性质}
\begin{cases}
\text{参数含义: } \mu \text{ 为位置参数(注意对称性),} \\
\sigma^2 \text{ 为形状参数} \\
\text{利用标准化求解概率}
\end{cases}
\end{cases}
\end{cases}
\end{cases} \\[1em]
\text{随机变量函数的概率分布}
\begin{cases}
\text{离散型: 关键取值及对应概率} \\
\text{连续型}
\begin{cases}
\text{定义法: 求解函数的分布函数} \\
\text{公式法: 适用单调的随机变量}
\end{cases}
\end{cases}
\end{cases}
\]

\begin{introduction}
  \item 随机变量~\ref{def:random-variable}
  \item 分布函数~\ref{def:distribution-function}
  \item 分布函数的性质~\ref{prop:distribution-function}
  \item 离散型随机变量~\ref{def:discrete-rv}
  \item 连续型随机变量~\ref{def:continuous-rv}
  \item 概率密度函数的性质~\ref{prop:pdf-property}
\end{introduction}

\section{随机变量与分布函数}

\subsection{随机变量}

\begin{definition}[随机变量] \label{def:random-variable}
设随机试验 \(E\) 的样本空间为 \(\Omega = \{\omega\}\)。如果对每一个 \(\omega \in \Omega\)，都有唯一的实数 \(X(\omega)\) 与之对应，并且对任意实数 \(x\)，集合
\[
\{\omega \mid X(\omega) \le x,\ \omega \in \Omega\}
\]
都是随机事件，则称定义在 \(\Omega\) 上的实值单值函数 \(X(\omega)\) 为\textbf{随机变量}，简记为 \(X\)。
\end{definition}

一般用大写字母 \(X, Y, Z, \dots\) 或希腊字母 \(\xi, \eta, \zeta, \dots\) 来表示随机变量。

\begin{note}
随机变量是"其值随试验结果而定"的变量，其定义有三个要点：
\begin{enumerate}
    \item 它是定义在样本空间 \(\Omega\) 上的实值单值函数——定义域是样本空间，而非实数集；
    \item 对任意实数 \(x\)，\(\{X \le x\}\) 必须是随机事件（即可定义概率）；
    \item 取值在试验之前不能预知，且取各值有一定的概率。
\end{enumerate}
随机事件是从\textbf{静态}的观点研究随机现象，而随机变量则引入了\textbf{动态}的观点（类似于高等数学中常量与变量的区别），使概率论的研究从对单个事件的孤立分析，上升到对随机变量整体规律的系统研究。
\end{note}

\begin{note}
随机变量一定范围内的取值本质上与随机事件等价。例如，"掷骰子出现的点数"是一个随机变量 \(X\)，则 \(X=3\)、\(X \le 4\)、\(X\) 为偶数等都是随机事件。
\end{note}

\subsection{分布函数}

\begin{definition}[分布函数] \label{def:distribution-function}
设 \(X\) 是随机变量，\(x\) 是任意实数，称函数
\[
F(x) = P\{X \le x\} \quad (-\infty < x < +\infty)
\]
为随机变量 \(X\) 的\textbf{分布函数}，也称 \(X\) 服从 \(F(x)\) 分布，记为 \(X \sim F(x)\)。
\end{definition}

当 \(x\) 从 \(-\infty\) 到 \(+\infty\) 取遍所有实数时，\(\{X \le x\}\) 就相应从不可能事件 \(\varnothing\) 取到必然事件 \(\Omega\)，因此分布函数\textbf{完整地描述了随机变量的概率规律}。


\subsection{分布函数的性质}

\begin{property}\label{prop:distribution-function}
设 \(F(x)\) 为随机变量 \(X\) 的分布函数，则它满足以下三条性质：
\begin{enumerate}
    \item \textbf{单调不减}：对任意实数 \(x_1 < x_2\)，有 \(F(x_1) \le F(x_2)\)；
    \item \textbf{有界性}：\(0 \le F(x) \le 1\)，且
    \[
    F(-\infty) = \lim_{x \to -\infty} F(x) = 0, \quad F(+\infty) = \lim_{x \to +\infty} F(x) = 1;
    \]
    \item \textbf{右连续}：对任意 \(x_0 \in \mathbf{R}\)，有 \(\lim_{x \to x_0^+} F(x) = F(x_0)\)。
\end{enumerate}
\end{property}

\begin{remark}
上述三条性质也是\textbf{充要条件}：若函数 \(F(x)\) 满足性质 (1)--(3)，则它一定是某个随机变量的分布函数。因此，这三条性质常用于判断一个函数能否作为分布函数。
\end{remark}


\subsection{分布函数与事件概率}

分布函数最重要的应用是通过它来计算各种事件的概率。记 \(F(a^-) = \lim_{x \to a^-} F(x)\) 为 \(F\) 在点 \(a\) 处的左极限值，则：

\begin{property}\label{prop:df-probability}
设随机变量 \(X\) 的分布函数为 \(F(x)\)，则：
\begin{enumerate}
    \item \(P\{X \le a\} = F(a)\)；
    \item \(P\{X < a\} = F(a^-)\)；
    \item \(P\{X = a\} = F(a) - F(a^-)\)；
    \item \(P\{a < X < b\} = F(b^-) - F(a)\)；
    \item \(P\{a < X \le b\} = F(b) - F(a)\)；
    \item \(P\{a \le X < b\} = F(b^-) - F(a^-)\)；
    \item \(P\{a \le X \le b\} = F(b) - F(a^-)\)；
    \item \(P\{X > a\} = 1 - F(a)\)。
\end{enumerate}
\end{property}

\begin{note}
记忆技巧：事件中包含等号 \(\le\) 或 \(=\) 的一端取函数值 \(F(\cdot)\)，不包含等号的一端（开区间端）取左极限 \(F(\cdot^-)\)。例如：
\begin{itemize}
    \item \(P\{a < X \le b\}\) 中，\(b\) 端有等号取 \(F(b)\)，\(a\) 端无等号取 \(F(a^-)\)；
    \item \(P\{X = a\} = F(a) - F(a^-)\) 即函数值减去左极限。
\end{itemize}
特别地，当 \(F(x)\) 在点 \(a\) 处连续时，\(F(a^-) = F(a)\)，从而 \(P\{X = a\} = 0\)。
\end{note}

\begin{exercise}\label{exer:2-1}
设随机变量 \(X\) 的分布函数为
\[
F(x) = \begin{cases}
0, & x < 0, \\
\dfrac{1}{2}, & 0 \le x < 1, \\
1 - e^{-x}, & x \ge 1.
\end{cases}
\]
求 \(P\{X = 0\}\)、\(P\{0 < X \le 2\}\) 和 \(P\{X > 1\}\)。
\end{exercise}

\begin{solution}
\begin{enumerate}
    \item \(P\{X = 0\} = F(0) - F(0^-) = \dfrac{1}{2} - 0 = \dfrac{1}{2}\)；
    \item \(P\{0 < X \le 2\} = F(2) - F(0) = (1 - e^{-2}) - \dfrac{1}{2} = \dfrac{1}{2} - e^{-2}\)；
    \item \(P\{X > 1\} = 1 - F(1) = 1 - (1 - e^{-1}) = e^{-1}\)。
\end{enumerate}
\end{solution}

\begin{exercise}\label{exer:2-2}
下列函数中，可以作为某个随机变量的分布函数的是（\quad）。\par
\begin{tabular}{ll}
A. \(F(x) = \dfrac{1}{\pi}\arctan x + \dfrac{1}{2},\ x \in \mathbf{R}\) & B. \(F(x) = \dfrac{1}{1+x^2},\ x \in \mathbf{R}\) \\[6pt]
C. \(F(x) = \begin{cases} 0, & x < 0 \\ 1-x, & 0 \le x < 1 \\ 1, & x \ge 1 \end{cases}\) & D. \(F(x) = \begin{cases} \sin x, & 0 \le x \le \pi \\ 0, & \text{其他} \end{cases}\)
\end{tabular}
\end{exercise}

\begin{solution}
应选 A。

逐条验证分布函数的三条充要条件：
\begin{itemize}
    \item A：\(\arctan x\) 单调递增，故 \(F(x)\) 单调不减；\(F(-\infty)=0\)，\(F(+\infty)=1\)；\(\arctan x\) 连续，故右连续。\textbf{三条均满足}。
    \item B：\(F(x)\) 在 \((-\infty,+\infty)\) 上\textbf{不是}单调不减的（例如 \(F(0)=1\)，\(F(1)=\frac{1}{2}\)），排除。
    \item C：在 \([0,1)\) 上 \(F(x)=1-x\) 单调\textbf{递减}，不满足单调不减，排除。
    \item D：\(\sin x\) 在 \([0,\pi]\) 上不是单调函数，不满足单调不减，排除。
\end{itemize}
\end{solution}

\begin{exercise}\label{exer:2-3}
设随机变量 \(X\) 的分布函数为
\[
F(x) = \begin{cases}
a + b e^{-x}, & x > 0, \\
0, & x \le 0.
\end{cases}
\]
求常数 \(a, b\) 及概率 \(P\{|X| < 2\}\)。
\end{exercise}

\begin{solution}
\(F(x)\) 中含有 2 个未知参数，依据分布函数的性质建立两个独立方程解之。

由分布函数的有界性，有
\[
F(+\infty) = \lim_{x \to +\infty} F(x) = \lim_{x \to +\infty} (a + b e^{-x}) = a = 1,
\]
得 \(a = 1\)。又由分布函数的右连续性，在 \(x=0\) 处有
\[
F(0) = \lim_{x \to 0^+} F(x) = a + b = 0,
\]
得 \(b = -1\)。所以
\[
F(x) = \begin{cases}
1 - e^{-x}, & x > 0, \\
0, & x \le 0.
\end{cases}
\]
从而
\[
P\{|X| < 2\} = P\{-2 < X < 2\} = F(2^-) - F(-2) = 1 - e^{-2} - 0 = 1 - e^{-2}.
\]
\end{solution}





















\section{颜色主题}

\begin{table}[htbp]
  \caption{ElegantBook 模板中的颜色主题\label{tab:color thm}}
  \centering
  \begin{tabular}{ccccccc}
  \toprule
    & \textcolor{structure1}{green} 
    & \textcolor{structure2}{cyan} 
    & \textcolor{structure3}{blue}
    & \textcolor{structure4}{gray} 
    & \textcolor{structure5}{black} 
    & 主要使用的环境\\
  \midrule
    structure & \ccr{structure1}
    & \ccr{structure2}
    & \ccr{structure3} 
    & \ccr{structure4} 
    & \ccr{structure5} 
    & chapter \ section \ subsection \\
    main      & \ccr{main1}
    & \ccr{main2}
    & \ccr{main3}
    & \ccr{main4}
    & \ccr{main5}
    & definition \ exercise \ problem \\
    second    & \ccr{second1}
    & \ccr{second2}
    & \ccr{second3}
    & \ccr{second4}
    & \ccr{second5}
    & theorem \ lemma \ corollary\\
    third     & \ccr{third1}
    & \ccr{third2}
    & \ccr{third3}
    & \ccr{third4}
    & \ccr{third5}
    & proposition\\
  \bottomrule
  \end{tabular}
\end{table}

如果需要自定义颜色的话请选择 \lstinline{nocolor} 选项或者使用 \lstinline{color=none}，然后在导言区定义 structurecolor、main、second、third 颜色，具体方法如下：
\begin{lstlisting}[tabsize=4]
\definecolor{structurecolor}{RGB}{0,0,0}
\definecolor{main}{RGB}{70,70,70}    
\definecolor{second}{RGB}{115,45,2}    
\definecolor{third}{RGB}{0,80,80}
\end{lstlisting}

\section{封面}

\subsection{封面个性化}

从 3.10 版本开始，封面更加弹性化，用户可以自行选择输出的内容，包括 \lstinline{\title} 在内的所有封面元素都可为空。目前封面的元素有

\begin{table}[htbp]
  \centering
  \caption{封面元素信息}
  \begin{tabular}{p{0.07\textwidth}p{0.15\textwidth}|p{0.07\textwidth}p{0.15\textwidth}|p{0.07\textwidth}p{0.15\textwidth}}
    \toprule
    信息 & 命令 & 信息 & 命令 & 信息 & 命令 \\
    \midrule
    标题 & \lstinline|\title| & 副标题 & \lstinline|\subtitle| & 作者 & \lstinline|\author| \\
    机构 & \lstinline|\institute| & 日期 &  \lstinline|\date| & 版本 & \lstinline|\version| \\
    箴言 & \lstinline|\extrainfo| & 封面图 & \lstinline|\cover| & 徽标 & \lstinline|\logo| \\
    \bottomrule
  \end{tabular}
\end{table}

另外，额外增加一个 \lstinline{\bioinfo} 命令，有两个选项，分别是信息标题以及信息内容。比如需要显示{\kaishu User Name：111520}，则可以使用 
\begin{lstlisting}
\bioinfo{User Name}{115520}
\end{lstlisting}

封面中间位置的色块的颜色可以使用下面命令进行修改：
\begin{lstlisting}
\definecolor{customcolor}{RGB}{32,178,170}
\colorlet{coverlinecolor}{customcolor}
\end{lstlisting}

\subsection{封面图}

本模板使用的封面图片来源于 \href{https://pixabay.com/en/tea-time-poetry-coffee-reading-3240766/}{pixabay.com}\footnote{感谢 China\TeX{} 提供免费图源网站，另外还推荐 \href{https://www.pexels.com/}{pexels.com}。}，图片完全免费，可用于任何场景。封面图片的尺寸为 $1280 \times 1024$, 更换图片的时候请\textbf{严格}按照封面图片尺寸进行裁剪。推荐一个免费的在线图片裁剪网站 \href{https://www.fotor.com/cn}{fotor.com}。用户 QQ 群内有一些合适尺寸的封面，欢迎取用。

\subsection{徽标}

本文用到的 Logo 比例为 1:1，也即正方形图片，在更换图片的时候请选择合适的图片进行替换。

\subsection{自定义封面}

另外，如果使用自定义的封面，比如 Adobe illustrator 或者其他软件制作的 A4 PDF 文档，请把 \lstinline{\maketitle} 注释掉，然后借助 \lstinline{pdfpages} 宏包将自制封面插入即可。如果使用 \lstinline{titlepage} 环境，也是类似。如果需要 2.x 版本的封面，请参考 \href{https://github.com/EthanDeng/etitlepage}{etitlepage}。

\section{章标标题}

本模板内置 2 套\textit{章标题显示风格}，包含 \lstinline{hang}（默认）与 \lstinline{display} 两种风格，区别在于章标题单行显示（\lstinline{hang}）与双行显示（\lstinline{display}），本说明使用了 \lstinline{hang}。调用方式为
\begin{lstlisting}
\documentclass[hang]{elegantbook} %or
\documentclass[titlestyle=hang]{elegantbook}
\end{lstlisting}

在章标题内，章节编号默认是以数字显示，也即{\kaishu 第 1 章}，{\kaishu 第 2 章}等等，如果想要把数字改为中文，可以使用
\begin{lstlisting}
\documentclass[chinese]{elegantbook} %or
\documentclass[scheme=chinese]{elegantbook}
\end{lstlisting}

\section{数学环境简介}

在我们这个模板中，我们定义了两种不同的定理模式 \lstinline{mode}，包括简单模式（\lstinline{simple}）和炫彩模式（\lstinline{fancy}），默认为 \lstinline{fancy} 模式，不同模式的选择为
\begin{lstlisting}
\documentclass[simple]{elegantbook} %or
\documentclass[mode=simple]{elegantbook}
\end{lstlisting}

本模板定义了四大类环境

\begin{itemize}
  \item \textit{定理类环境}，包含标题和内容两部分，全部定理类环境的编号均以章节编号。根据格式的不同分为 3 种
    \begin{itemize}
      \item \textcolor{main}{\textbf{definition}} 环境，颜色为 \textcolor{main}{main}；
      \item \textcolor{second}{\textbf{theorem、lemma、corollary、axiom、postulate}} 环境，颜色为 \textcolor{second} {second}；
      \item \textcolor{third}{\textbf{proposition}} 环境，颜色为 \textcolor{third}{third}。
    \end{itemize}
  \item \textit{示例类环境}，有 \textbf{example、problem、exercise} 环境（对应于例、例题、练习），自动编号，编号以章节为单位，其中 \textbf{exercise} 有提示符。
  \item \textit{提示类环境}，有 \textbf{note} 环境，特点是：无编号，有引导符。
  \item \textit{结论类环境}，有 \textbf{conclusion、assumption、property、remark、solution} 环境\footnote{本模板还添加了一个 result 选项，用于隐藏 \lstinline{solution} 和 \lstinline{proof} 环境，默认为显示（\lstinline{result=answer}），隐藏使用 \lstinline{result=noanswer}。}，三者均以粗体的引导词为开头，和普通段落格式一致。
\end{itemize}

\subsection{定理类环境的使用}

由于本模板使用了 \lstinline{tcolorbox} 宏包来定制定理类环境，所以和普通的定理环境的使用有些许区别，定理的使用方法如下：
\begin{lstlisting}
\begin{theorem}{theorem name}{label}
  The content of theorem.
\end{theorem}
\end{lstlisting}

第一个必选项 \lstinline{theorem name} 是定理的名字，第二个必选项 \lstinline{label} 是交叉引用时所用到的标签，交叉引用的方法为 \verb|\ref{thm:label}|。请注意，交叉引用时必须加上前缀 \lstinline{thm:}。

在用户多次反馈下，4.x 之后，引入了原生定理的支持方式，也就是使用可选项方式：

\begin{lstlisting}
\begin{theorem}[theorem name] \label{thm:theorem-label}
  The content of theorem.
\end{theorem}
% or 
\begin{theorem} \label{thm:theorem-withou-name}
  The content of theorem without name.
\end{theorem}
\end{lstlisting}

其他相同用法的定理类环境有：

\begin{table}[htbp]
   \centering
   \caption{定理类环境}
     \begin{tabular}{llll}
     \toprule
     环境名 & 标签名 & 前缀 & 交叉引用 \\
     \midrule
     definition & label & def   & \lstinline|\ref{def:label}| \\
     theorem & label & thm   & \lstinline|\ref{thm:label}| \\
     postulate & label & pos & \lstinline|\ref{pos:label}| \\
     axiom & label & axi & \lstinline|\ref{axi:label}|\\
     lemma & label & lem   & \lstinline|\ref{lem:label}| \\
     corollary & label & cor   & \lstinline|\ref{cor:label}| \\
     proposition & label & pro   & \lstinline|\ref{pro:label}| \\
     \bottomrule
     \end{tabular}%
   \label{tab:theorem-class}%
 \end{table}%
 
% \subsection{算法环境}

 
% \begin{algorithm}\label{alg:test}
%   \Input{A bitmap $I$ of size $w \times l$}
%   \Output{A partition of the bitmap}
%   \BlankLine
%   \emph{special treatment of the first line}\;
%   \For{$i \leftarrow 2$ \KwTo $l$}{
%     \emph{special treatment of the first element of line $i$}\;
%     \For{$j \leftarrow 2$ \KwTo $w$}{\label{forins}
%       $\Left \leftarrow \FindCompress{$I[i,j-1]$}$\;
%       $\Up \leftarrow \FindCompress{$I[i-1,]$}$\;
%       $\This \leftarrow \FindCompress{$I[i,j]$}$\;
%       \If(\tcp*[h]{O(\Left,\This)==1}){\Left compatible with \This}{\label{lt}
%         \lIf{$\Left < \This$}{$\Union{\Left,\This}$}
%         \lElse{$\Union{\This,\Left}$}
%       }
%       \If(\tcp*[f]{O(\Up,\This)==1}){\Up compatible with \This}{\label{ut}
%         \lIf{$\Up < \This$}{$\Union{\Up,\This}$}
%         \tcp{\This is put under \Up to keep tree as flat as possible}\label{cmt}
%         \lElse{$\Union{\This,\Up}$}\tcp*[r]{\This{} linked to \Up}\label{lelse}
%       }
%     }
%     \lForEach{element $e$ of the line $i$}{\FindCompress{p}}
%   }
%   \caption{disjoint decomposition}\label{algo_disjdecomp}
% \end{algorithm}


\subsection{修改计数器}

当前定理等环境计数器按章计数，如果想修改定理类环境按节计数，可以修改计数器选项 thmcnt：

\begin{lstlisting}
  \documentclass[section]{elegantbook} %or
  \documentclass[thmcnt=section]{elegantbook}
\end{lstlisting}


\subsection{其他环境的使用}

其他三种环境没有选项，可以直接使用，比如 \lstinline{example} 环境的使用方法与效果：
\begin{lstlisting}
\begin{example}
   This is the content of example environment.
\end{example}
\end{lstlisting}

这几个都是同一类环境，区别在于

\begin{itemize}
  \item 示例环境（example）、练习（exercise）与例题（problem）章节自动编号；
  \item 注意（note），练习（exercise）环境有提醒引导符；
  \item 结论（conclusion）等环境都是普通段落环境，引导词加粗。
\end{itemize}

\section{列表环境}
本模板借助于 \lstinline{tikz} 定制了 \lstinline{itemize} 和 \lstinline{enumerate} 环境，其中 \lstinline{itemize} 环境修改了 3 层嵌套，而 \lstinline{enumerate} 环境修改了 4 层嵌套（仅改变颜色）。示例如下\\[2ex]
\begin{minipage}[b]{0.49\textwidth}
  \begin{itemize}
    \item first item of nesti;
    \item second item of nesti;
      \begin{itemize}
        \item first item of nestii;
        \item second item of nestii;
        \begin{itemize}
          \item first item of nestiii;
          \item second item of nestiii.
        \end{itemize}   
      \end{itemize}
  \end{itemize}
\end{minipage}
\begin{minipage}[b]{0.49\textwidth}
  \begin{enumerate}
    \item first item of nesti;
    \item second item of nesti;
      \begin{enumerate}
        \item first item of nestii;
        \item second item of nestii;
        \begin{enumerate}
          \item first item of nestiii;
          \item second item of nestiii.
        \end{enumerate}   
      \end{enumerate}
  \end{enumerate}
\end{minipage}

\section{参考文献}
文献部分，本模板调用了 biblatex 宏包，并提供了 biber（默认） 和 bibtex 两个后端选项，可以使用 \lstinline{bibend} 进行修改：

\begin{lstlisting}
  \documentclass[bibtex]{elegantbook}
  \documentclass[bibend=bibtex]{elegantbook}
\end{lstlisting}

关于文献条目（bib item），你可以在谷歌学术，Mendeley，Endnote 中取，然后把它们添加到 \lstinline{reference.bib} 中。在文中引用的时候，引用它们的键值（bib key）即可。

为了方便文献样式修改，模板引入了 \lstinline{bibstyle} 和 \lstinline{citestyle} 选项，默认均为数字格式（numeric），参考文献示例：\cite{cn1,en2,en3} 使用了中国一个大型的 P2P 平台（人人贷）的数据来检验男性投资者和女性投资者在投资表现上是否有显著差异。

如果需要设置为国标 GB7714-2015，需要使用：
\begin{lstlisting}
  \documentclass[citestyle=gb7714-2015, bibstyle=gb7714-2015]{elegantbook} 
\end{lstlisting}

如果需要添加排序方式，可以在导言区加入
\begin{lstlisting}
  \ExecuteBibliographyOptions{sorting=ynt}
\end{lstlisting}

启用国标之后，可以加入 \lstinline{sorting=gb7714-2015}。

\section{添加序章}

如果你想在第一章前面添序章，不改变原本章节序号，可以在第一章内容前面使用 
\begin{lstlisting}
\chapter*{Introduction}
\markboth{Introduction}{Introduction}
The content of introduction.
\end{lstlisting}

\section{目录选项与深度}
本模板添加了一个目录选项 \lstinline{toc}，可以设置目录为单栏（\lstinline{onecol}）和双栏（\lstinline{twocol}）显示，比如双栏显示可以使用
\begin{lstlisting}
\documentclass[twocol]{elegantbook}
\documentclass[toc=twocol]{elegantbook}
\end{lstlisting}

默认本模板目录深度为 1，你可以在导言区使用
\begin{lstlisting}
\setcounter{tocdepth}{2}
\end{lstlisting}
将其修改为 2 级目录（章与节）显示。


\section{章节摘要}
模板新增了一个章节摘要环境（introduction），使用示例
\begin{lstlisting}
\begin{introduction}
  \item Definition of Theorem
  \item Ask for help
  \item Optimization Problem
  \item Property of Cauchy Series
  \item Angle of Corner
\end{introduction}
\end{lstlisting}
效果如下：
\begin{introduction}
  \item Definition of Theorem
  \item Ask for help
  \item Optimization Problem
  \item Property of Cauchy Series
  \item Angle of Corner
\end{introduction}

环境的标题文字可以通过这个环境的可选参数进行修改，修改方法为：
\begin{lstlisting}
\begin{introduction}[Brief Introduction]
...
\end{introduction}
\end{lstlisting}

\section{章后习题}
前面我们介绍了例题和练习两个环境，这里我们再加一个，章后习题（\lstinline{problemset}）环境，用于在每一章结尾，显示本章的练习。使用方法如下

\begin{lstlisting}
\begin{problemset}
  \item exercise 1
  \item exercise 2
  \item exercise 3
\end{problemset}
\end{lstlisting}


效果如下：
\begin{problemset}
  \item exercise 1
  \item exercise 2
  \item exercise 3
  \item 测试数学公式
  \begin{equation}
    a^2+b^2=c_{2_{i}} (1,2) [1,23]
  \end{equation}
\end{problemset}

\begin{remark}
如果你想把 \lstinline{problemset} 环境的标题改为其他文字，你可以类似于 introduction 环境修改 problemset 的可选参数。另外，目前这个环境会自动出现在目录中，但是不会出现在页眉页脚信息中（待解决）。
\end{remark}

\begin{solution}
如果你想把 \lstinline{problemset} 环境的标题改为其他文字，你可以类似于 introduction 环境修改 problemset 的可选参数。另外，目前这个环境会自动出现在目录中，但是不会出现在页眉页脚信息中（待解决）。
\end{solution}





\section{旁注}

在 3.08 版本中，我们引入了 旁注设置选项 \lstinline{marginpar=margintrue} 以及测试命令 \lstinline{\elegantpar} ，但是由此带来一堆问题。我们决定在 3.09 版本中将其删除，并且，在旁注命令得到大幅度优化之前，不会将此命令再次引入书籍模板中。对此造成各位用户的不方便，非常抱歉！不过我们保留了 \lstinline{marginpar} 这个选项，你可以使用 \lstinline{marginpar=margintrue} 获得保留右侧旁注的版面设计。然后使用系统自带的 \lstinline{\marginpar} 或者 \lstinline{marginnote} 宏包的 \lstinline{\marginnote} 命令。

\begin{remark}
在使用旁注的时候，需要注意的是，文本和公式可以直接在旁注中使用。

\begin{lstlisting}
% text
\marginpar{margin paragraph text}

% equation
\marginpar{
  \begin{equation}
    a^2 + b^2 = c^2
  \end{equation}
}
\end{lstlisting}

但是浮动体（表格、图片）需要注意，不能用浮动体环境，需要使用直接插图命令或者表格命令环境。然后使用 \lstinline{\captionof} 为其设置标题。为了得到居中的图表，可以使用 \lstinline{\centerline} 命令或者 \lstinline{center} 环境。更多详情请参考：\href{https://tex.stackexchange.com/questions/5583/caption-of-figure-in-marginpar-and-caption-of-wrapfigure-in-margin}{Caption of Figure in Marginpar}。

\begin{lstlisting}
% graph with centerline command
\marginpar{
  \centerline{
    \includegraphics[width=0.2\textwidth]{logo.png}
  }
  \captionof{figure}{your figure caption}
}

% graph with center environment
\marginpar{
  \begin{center}
    \includegraphics[width=0.2\textwidth]{logo.png}
    \captionof{figure}{your figure caption}
  \end{center}
}
\end{lstlisting}

\end{remark}

\chapter{字体选项}
字体选项独立成章的原因是，我们希望本模板的用户关心模板使用的字体，知晓自己使用的字体以及遇到字体相关的问题能更加便捷地找到答案。

\textcolor{red}{\bfseries 重要提示}：从 3.10 版本更新之后，沿用至今的 newtx 系列字体被重新更改为 cm 字体。并且新增中文字体（\lstinline{chinesefont}）选项。

\section{数学字体选项}

本模板定义了一个数学字体选项（\lstinline{math}），可选项有三个：
\begin{enumerate}
  \item \lstinline{math=cm}（默认），使用 \LaTeX{} 默认数学字体（推荐，无需声明）；
  \item \lstinline{math=newtx}，使用 \lstinline{newtxmath} 设置数学字体（潜在问题比较多）。
  \item \lstinline{math=mtpro2}，使用 \lstinline{mtpro2} 宏包设置数学字体，要求用户已经成功安装此宏包。
\end{enumerate}

\section{使用 newtx 系列字体}

如果需要使用原先版本的 \lstinline{newtx} 系列字体，可以通过显示声明数学字体：

\begin{lstlisting}
\documentclass[math=newtx]{elegantbook}
\end{lstlisting}

\subsection{连字符}

如果使用 \lstinline{newtx} 系列字体宏包，需要注意下连字符的问题。
\begin{equation}
  \int_{R^q} f(x,y) dy.\emph{of\kern0pt f}
\end{equation}
的代码为
\begin{lstlisting}
\begin{equation}
  \int_{R^q} f(x,y) dy.\emph{of \kern0pt f}
\end{equation}
\end{lstlisting}

\subsection{宏包冲突}

另外在 3.08 版本中，有用户反馈模板在和 \lstinline{yhmath} 以及 \lstinline{esvect} 等宏包搭配使用的时候会出现报错：
\begin{lstlisting}
LaTeX Error:
   Too many symbol fonts declared.
\end{lstlisting}

原因是在使用 \lstinline{newtxmath} 宏包时，重新定义了数学字体用于大型操作符，达到了 {\heiti 最多 16 个数学字体} 的上限，在调用其他宏包的时候，无法新增数学字体。为了减少调用非常用宏包，在此给出如何调用 \lstinline{yhmath} 以及 \lstinline{esvect} 宏包的方法。

请在 \lstinline{elegantbook.cls} 内搜索 \lstinline{yhmath} 或者 \lstinline{esvect}，将你所需要的宏包加载语句\textit{取消注释}即可。
\begin{lstlisting}
%%% use yhmath pkg, uncomment following code
% \let\oldwidering\widering
% \let\widering\undefined
% \RequirePackage{yhmath}
% \let\widering\oldwidering

%%% use esvect pkg, uncomment following code
% \RequirePackage{esvect}
\end{lstlisting}

\section{中文字体选项}
模板从 3.10 版本提供中文字体选项 \lstinline{chinesefont}，可选项有
\begin{enumerate}
\item \lstinline{ctexfont}：默认选项，使用 \lstinline{ctex} 宏包根据系统自行选择字体，可能存在字体缺失的问题，更多内容参考 \lstinline{ctex} 宏包\href{https://ctan.org/pkg/ctex}{官方文档}\footnote{可以使用命令提示符，输入 \lstinline{texdoc ctex} 调出本地 \lstinline{ctex} 宏包文档}。
\item \lstinline{founder}：方正字体选项（\textbf{需要安装方正字体}），后台调用 \lstinline{ctex} 宏包并且使用 \lstinline{fontset=none} 选项，然后设置字体为方正四款免费字体，方正字体下载注意事项见后文，用户只需要安装方正字体即可使用该选项。
\item \lstinline{nofont}：后台会调用 \lstinline{ctex} 宏包并且使用 \lstinline{fontset=none} 选项，不设定中文字体，用户可以自行设置中文字体，具体见后文。
\end{enumerate}

\subsection{方正字体选项}
由于使用 \lstinline{ctex} 宏包默认调用系统已有的字体，部分系统字体缺失严重，因此，用户希望能够使用其它字体，我们推荐使用方正字体。方正的{\songti 方正书宋}、{\heiti 方正黑体}、{\kaishu 方正楷体}、{\fangsong 方正仿宋}四款字体均可免费试用，且可用于商业用途。用户可以自行从\href{http://www.foundertype.com/}{方正字体官网}下载此四款字体，在下载的时候请\textbf{务必}注意选择 GBK 字符集，也可以使用 \href{https://www.latexstudio.net/}{\LaTeX{} 工作室}提供的\href{https://pan.baidu.com/s/1BgbQM7LoinY7m8yeP25Y7Q}{方正字体，提取码为：njy9} 进行安装。安装时，{\kaishu Win 10 用户请右键选择为全部用户安装，否则会找不到字体。}

\begin{figure}[!htb]
\centering
\includegraphics[width=0.9\textwidth]{founder.png}
\end{figure}

\subsection{其他中文字体}
如果你想完全自定义字体\footnote{这里仍然以方正字体为例。}，你可以选择 \lstinline{chinesefont=nofont}，然后在导言区设置
\begin{lstlisting}
\setCJKmainfont[BoldFont={FZHei-B01},ItalicFont={FZKai-Z03}]{FZShuSong-Z01}
\setCJKsansfont[BoldFont={FZHei-B01}]{FZKai-Z03}
\setCJKmonofont[BoldFont={FZHei-B01}]{FZFangSong-Z02}
\setCJKfamilyfont{zhsong}{FZShuSong-Z01}
\setCJKfamilyfont{zhhei}{FZHei-B01}
\setCJKfamilyfont{zhkai}[BoldFont={FZHei-B01}]{FZKai-Z03}
\setCJKfamilyfont{zhfs}[BoldFont={FZHei-B01}]{FZFangSong-Z02}
\newcommand*{\songti}{\CJKfamily{zhsong}}
\newcommand*{\heiti}{\CJKfamily{zhhei}}
\newcommand*{\kaishu}{\CJKfamily{zhkai}}
\newcommand*{\fangsong}{\CJKfamily{zhfs}}
\end{lstlisting}

\chapter{ElegantBook 写作示例}

\begin{introduction}
  \item 随机试验与样本空间~\ref{def:random-experiment}
  \item 随机事件及其运算~\ref{def:random-event}
  \item 完备事件组~\ref{def:complete-event-group}
  \item 概率的公理化定义~\ref{def:probability-axiom}
  \item 概率的基本性质~\ref{prop:probability-basic}
  \item 古典概型~\ref{def:classical-prob}
  \item 几何概型~\ref{def:geometric-prob}
  \item 条件概率~\ref{def:conditional-prob}
  \item 乘法公式~\ref{thm:multiplication}
  \item 全概率公式与贝叶斯公式~\ref{thm:total-probability}
  \item 事件的独立性~\ref{def:independence}
  \item 伯努利概型~\ref{def:bernoulli}
\end{introduction}

\section{Lebesgue 积分}
在前面各章做了必要的准备后，本章开始介绍新的积分。在 Lebesgue 测度理论的基础上建立了 Lebesgue 积分，其被积函数和积分域更一般，可以对有界函数和无界函数统一处理。正是由于 Lebesgue 积分的这些特点，使得 Lebesgue 积分比 Riemann 积分具有在更一般条件下的极限定理和累次积分交换积分顺序的定理，这使得 Lebesgue 积分不仅在理论上更完善，而且在计算上更灵活有效。

Lebesgue 积分有几种不同的定义方式。我们将采用逐步定义非负简单函数，非负可测函数和一般可测函数积分的方式。

由于现代数学的许多分支如概率论、泛函分析、调和分析等常常用到一般空间上的测度与积分理论，在本章最后一节将介绍一般的测度空间上的积分。

\subsection{积分的定义}

我们将通过三个步骤定义可测函数的积分。首先定义非负简单函数的积分。以下设 $E$ 是 $\mathcal{R}^n$ 中的可测集。

\begin{definition}[可积性] \label{def:int} 
设 $ f(x)=\sum\limits_{i=1}^{k} a_i \chi_{A_i}(x)$ 是 $E$ 上的\textbf{非负简单函数}，中文其中 $\{A_1,A_2,\ldots,A_k\}$ 是 $E$ 上的一个可测分割，$a_1,a_2,\ldots,a_k$ 是非负实数。定义 $f$ 在 $E$ 上的积分为 $\int_{a}^b f(x)$
\begin{equation}
   \label{inter}
   \int_{E} f dx = \sum_{i=1}^k a_i m(A_i) \pi \alpha\beta\sigma\gamma\nu\xi\epsilon\varepsilon. \oint_{a}^b\ointop_{a}^b\prod_{i=1}^n
\end{equation}
一般情况下 $0 \leq \int_{E} f dx \leq \infty$。若 $\int_{E} f dx < \infty$，则称 $f$ 在 $E$ 上可积。
\end{definition}

一个自然的问题是，Lebesgue 积分与我们所熟悉的 Riemann 积分有什么联系和区别？在 4.4 在我们将详细讨论 Riemann 积分与 Lebesgue 积分的关系。这里只看一个简单的例子。设 $D(x)$ 是区间 $[0,1]$ 上的 Dirichlet 函数。即 $D(x)=\chi_{Q_0}(x)$，其中 $Q_0$ 表示 $[0,1]$ 中的有理数的全体。根据非负简单函数积分的定义，$D(x)$ 在 $[0,1]$ 上的 Lebesgue 积分为
\begin{equation}
   \label{inter2}
   \int_0^1 D(x)dx = \int_0^1 \chi_{Q_0} (x) dx = m(Q_0) = 0
\end{equation}
即 $D(x)$ 在 $[0,1]$ 上是 Lebesgue 可积的并且积分值为零。但 $D(x)$ 在 $[0,1]$ 上不是 Riemann 可积的。


有界变差函数是与单调函数有密切联系的一类函数。有界变差函数可以表示为两个单调递增函数之差。与单调函数一样，有界变差函数几乎处处可导。与单调函数不同，有界变差函数类对线性运算是封闭的，它们构成一线空间。练习题 \ref{exer:43} 是一个性质的证明。

\begin{exercise}\label{exer:43}
设 $f \notin\in L(\mathcal{R}^1)$，$g$ 是 $\mathcal{R}^1$ 上的有界可测函数。证明函数
\begin{equation}
   \label{ex:1}
   I(t) = \int_{\mathcal{R}^1} f(x+t)g(x)dx \quad t \in \mathcal{R}^1
\end{equation}
是 $\mathcal{R}^1$ 上的连续函数。 
\end{exercise}

\begin{solution}
即 $D(x)$ 在 $[0,1]$ 上是 Lebesgue 可积的并且积分值为零。但 $D(x)$ 在 $[0,1]$ 上不是 Riemann 可积的。
\end{solution}





\begin{proof}
即 $D(x)$ 在 $[0,1]$ 上是 Lebesgue 可积的并且积分值为零。但 $D(x)$ 在 $[0,1]$ 上不是 Riemann 可积的。
\end{proof}

\begin{theorem}[Fubini 定理] \label{thm:fubi} 
（1）若 $f(x,y)$ 是 $\mathcal{R}^p\times\mathcal{R}^q$ 上的非负可测函数，则对几乎处处的 $x\in \mathcal{R}^p$，$f(x,y)$ 作为 $y$ 的函数是 $\mathcal{R}^q$ 上的非负可测函数，$g(x)=\int_{\mathcal{R}^q}f(x,y) dy$ 是 $\mathcal{R}^p$ 上的非负可测函数。并且
\begin{equation}
   \label{eq:461}
   \int_{\mathcal{R}^p\times\mathcal{R}^q} f(x,y) dxdy=\int_{\mathcal{R}^p}\left(\int_{\mathcal{R}^q}f(x,y)dy\right)dx.
\end{equation}

（2）若 $f(x,y)$ 是 $\mathcal{R}^p\times\mathcal{R}^q$ 上的可积函数，则对几乎处处的 $x\in\mathcal{R}^p$，$f(x,y)$ 作为 $y$ 的函数是 $\mathcal{R}^q$ 上的可积函数，并且 $g(x)=\int_{\mathcal{R}^q}f(x,y) dy$ 是 $\mathcal{R}^p$ 上的可积函数。而且~\ref{eq:461} 成立。
\end{theorem}

\ref{thm:fubi}

\begin{note}
在本模板中，引理（lemma），推论（corollary）的样式和定理~\ref{thm:fubi} 的样式一致，包括颜色，仅仅只有计数器的设置不一样。
\end{note}

\begin{exercise}\label{exer:1-3}
设 \(A, B\) 为两个随机事件，\(P(A) = 0.4\)，\(P(A \cup B) = 0.7\)。\par
(1) 求 \(P(B - A)\)。\par
(2) 若事件 \(A\) 和事件 \(B\) 互不相容，求 \(P(B)\)。
\end{exercise}

\begin{solution}
(1) 由加法公式
\[
P(A \cup B) = P(A) + P(B) - P(AB) = P(A) + P(B - A),
\]
故
\[
P(B - A) = P(A \cup B) - P(A) = 0.7 - 0.4 = 0.3.
\]

(2) 若 \(A\) 与 \(B\) 互不相容，则 \(P(AB) = 0\)，此时加法公式简化为
\[
P(A \cup B) = P(A) + P(B),
\]
故
\[
P(B) = P(A \cup B) - P(A) = 0.7 - 0.4 = 0.3.
\]
\end{solution}


\datechange{2022/03/08}{版本 4.2 正式发布。}

\begin{change}
  \item 对于 newtx 系列宏包更新导致的字体 bug 的修复；
  \item 修缮目录格式，为了达到这个目的，重新改写 \lstinline{\chaptername} 的重定义语句；
  \item 增加日语 \lstinline{lang=jp} 设定。
  \item 这个版本为一个临时性版本，在 \TeX Live 2022 发布之后，将尽快发布 4.3 版本，由于对于中文的改动比较大，可能会出现预期之外的 bug，有问题可以在 QQ 群或者 Github 反馈。
\end{change}

\datechange{2018/12/16}{版本 3.01 正式发布}

\begin{change}
  \item 调整 \lstinline{ctex} 宏包。
  \item 说明文档增加更新内容。
\end{change}
\nocite{*}
\printbibliography[heading=bibintoc, title=\ebibname]
\appendix



\end{document}
